% Provenance: copied from the design workspace CambridgePQ/recursion.fast, `report/sections/A-lean-contract.tex`, on 2026-10-19. Read-only copy; the source is authoritative.
% A-lean-contract -- written by W7; listings regenerated from contract-v03 by scripts/regen-appendix-a.py (W10b)
% The group below adds a glyph fallback (e.g. U+2113) to the monospace font of the listings only.
\begingroup
\directlua{luaotfload.add_fallback("rfleanmono", {"DejaVu Sans:mode=node;", "STIX Two Math:mode=node;"})}
\setmonofont{DejaVu Sans Mono}[Scale=0.82,RawFeature={fallback=rfleanmono}]
\section{Lean contract listings}\label{sec:appendix-contract}

This appendix reproduces the trusted contract v0.3 (\code{contract-v03/RecursionFast/Contract/}),
which compiles against Lean v4.34.1 and Mathlib v4.34.1 and imports nothing else except the vendored
riscv-zkvm model. The files appear in the import order of the root module, which is the intended
reading order. Each listing is the \emph{complete} compiled source, generated from it by a script
(\path{report/scripts/regen-appendix-a.py}, which also has a \code{--check} mode for CI), with
three mechanical changes: the module header comment, paraphrased in the paragraph before the
listing, is omitted; the lines \lean{namespace RecursionFast} and \lean{end RecursionFast} that
enclose every file are omitted; and lines longer than 112 characters are broken at a space, with
the continuation indented four further spaces. Nothing is elided. The listings in
\cref{sec:protocol,sec:security} are excerpts of these. For each file the paragraph states what it
defines and what an auditor must check. \Cref{tab:app-loc} gives the line counts, and
\cref{subsec:app-checklist} the auditor's checklist.

\begin{table}[h]
\centering\small
\begin{tabular}{@{}lrrrrl@{}}
\toprule
File & Code & Comment & Blank & Total & v0.2 code \\
\midrule
\code{Field.lean} & 36 & 27 & 20 & 83 & 36 \\
\code{R1CS.lean} & 36 & 31 & 18 & 85 & 41 \\
\code{Shapes.lean} & 50 & 37 & 25 & 112 & 50 \\
\code{Security.lean} & 55 & 38 & 20 & 113 & 55 \\
\code{Protocol.lean} & 27 & 23 & 10 & 60 & 27 \\
\code{Timer.lean} (new) & 22 & 23 & 14 & 59 & -- \\
\code{Machine.lean} & 297 & 117 & 59 & 473 & 165 \\
\code{Agrees.lean} & 18 & 17 & 11 & 46 & 18 \\
\code{Cost.lean} & 73 & 61 & 35 & 169 & 38 \\
\code{QProg.lean} & 21 & 13 & 8 & 42 & 21 \\
\code{Accumulation.lean} & 77 & 48 & 29 & 154 & 72 \\
\code{Certificate.lean} & 72 & 35 & 19 & 126 & 48 \\
\code{Contract.lean} (root) & 12 & 1 & 1 & 14 & 11 \\
\midrule
\textbf{Total} & \textbf{796} & \textbf{471} & \textbf{269} & \textbf{1{,}536} & \textbf{582} \\
\bottomrule
\end{tabular}
\caption{Line counts of the trusted contract after the round-3 fixes (\code{contract-v03/scripts/loc.py},
re-run for this appendix; a line with code and a comment counts as code; doc comments count as comments). Not counted: the vendored riscv-zkvm model, of
which the contract reaches 474 declarations in 930 source lines (205 of them theorems pulled in by
\lean{deriving DecidableEq} on \lean{Instruction}).}
\label{tab:app-loc}
\end{table}

\subsection{Root module}\label{subsec:app-root}

The root module only imports the twelve files. An auditor checks that its import list is the
reading order below and that no file imports anything outside Lean core, Mathlib and
\code{RiscvZkvm} (the closure of \lean{import RecursionFast.Contract} has 2,635 modules and no
ArkLib, CompPoly or VCVio module).

\begin{leancode}[fontsize=\scriptsize,label={\footnotesize RecursionFast/Contract.lean}]
import RecursionFast.Contract.Field
import RecursionFast.Contract.R1CS
import RecursionFast.Contract.Shapes
import RecursionFast.Contract.Security
import RecursionFast.Contract.Protocol
import RecursionFast.Contract.Timer
import RecursionFast.Contract.Machine
import RecursionFast.Contract.Agrees
import RecursionFast.Contract.Cost
import RecursionFast.Contract.QProg
import RecursionFast.Contract.Accumulation
import RecursionFast.Contract.Certificate

/-! # recursion.fast trusted contract v0.3 (root). The imports above are the reading order. -/
\end{leancode}

\subsection{\texttt{Field.lean}: the fields \texorpdfstring{$K$}{K} and \texorpdfstring{$E$}{E}}

$K=\mathbb{F}_2[x]/(x^{64}+x^4+x^3+x+1)$ in the polynomial basis (bit $i$ of a word is the
coefficient of $x^i$), so every 64-bit word is an element and every 8-byte string a canonical
encoding; $E=K[y]/(y^3+y+1)$ as three limbs. Lean core only. The field axioms, $|K|=2^{64}$,
$|E|=2^{192}$, irreducibility of both moduli and the definitional identification with CompPoly's
\lean{BF64} and \lean{BF64.Ext3} are theorems of the untrusted library. \emph{Check}: the reduction
constant \lean{0x1B} is $x^4+x^3+x+1$; two folds suffice (after one fold the degree is below 67);
the $E$ product folds $y^3=y+1$ and $y^4=y^2+y$; the encodings are little-endian, limbs in the order
$c_0,c_1,c_2$; \lean{K.clmul} is also the semantics of Zbc.

\begin{leancode}[fontsize=\scriptsize,label={\footnotesize Contract/Field.lean}]
/-- A 64-bit machine word (also the unit of the tape and of oracle leaves). -/
abbrev Word := BitVec 64

/-- An element of `K = GF(2^64)`: bit `i` of `bits` is the coefficient of `x^i`. -/
structure K where
  bits : Word
  deriving DecidableEq

namespace K

/-- Carry-less product of two polynomials of degree `< 64` over GF(2), as a 128-bit word
(schoolbook shift-and-xor). Also the semantics of the Zbc instructions (`Machine.lean`). -/
def clmul (a b : Word) : BitVec 128 :=
  Fin.foldl 64 (fun acc i => if a.getLsbD i then acc ^^^ (b.zeroExtend 128 <<< (i : Nat)) else acc) 0

/-- One reduction fold `hi·x^64 + lo ↦ hi·(x^4 + x^3 + x + 1) + lo` (`x^64 ≡ 0x1B`). -/
def fold (t : BitVec 128) : BitVec 128 :=
  clmul ((t >>> 64).setWidth 64) 0x1B ^^^ (t.setWidth 64).zeroExtend 128

/-- Reduction modulo `x^64 + x^4 + x^3 + x + 1`: two folds bring the degree below 64. -/
def reduce (t : BitVec 128) : Word := (fold (fold t)).setWidth 64

instance : Zero K := ⟨⟨0⟩⟩
instance : One K := ⟨⟨1⟩⟩
/-- Addition = subtraction = xor (characteristic 2); negation is the identity. -/
instance : Add K := ⟨fun a b => ⟨a.bits ^^^ b.bits⟩⟩
instance : Sub K := ⟨fun a b => ⟨a.bits ^^^ b.bits⟩⟩
instance : Neg K := ⟨fun a => a⟩
/-- Multiplication: carry-less product, then reduction. -/
instance : Mul K := ⟨fun a b => ⟨reduce (clmul a.bits b.bits)⟩⟩

/-- Canonical encoding: the 8 little-endian bytes of `bits` (Rust `u64::to_le_bytes`). On the
machine an element is one doubleword, which memory stores little-endian, i.e. these bytes. -/
def toBytes (a : K) : List (BitVec 8) := List.ofFn fun i : Fin 8 => (a.bits >>> (8 * i.val)).setWidth 8

end K

/-- An element `c0 + c1·y + c2·y²` of `E = K[y]/(y^3 + y + 1) = GF(2^192)`. -/
structure E where
  c0 : K
  c1 : K
  c2 : K
  deriving DecidableEq

namespace E

instance : Zero E := ⟨⟨0, 0, 0⟩⟩
instance : One E := ⟨⟨1, 0, 0⟩⟩
instance : Add E := ⟨fun a b => ⟨a.c0 + b.c0, a.c1 + b.c1, a.c2 + b.c2⟩⟩
/-- Schoolbook product folded by `y^3 = y + 1`, `y^4 = y^2 + y` (9 products in K). -/
instance : Mul E := ⟨fun a b =>
  let p3 := a.c1 * b.c2 + a.c2 * b.c1
  let p4 := a.c2 * b.c2
  ⟨a.c0 * b.c0 + p3, a.c0 * b.c1 + a.c1 * b.c0 + p3 + p4, a.c0 * b.c2 + a.c1 * b.c1 + a.c2 * b.c0 + p4⟩⟩

/-- Canonical encoding as words: limbs `c0, c1, c2` in this order (24 bytes, each limb LE). -/
def toWords (a : E) : List Word := [a.c0.bits, a.c1.bits, a.c2.bits]

/-- Canonical 24-byte encoding: limbs `c0, c1, c2`, each 8 bytes little-endian. -/
def toBytes (a : E) : List (BitVec 8) := a.c0.toBytes ++ a.c1.toBytes ++ a.c2.toBytes

end E
\end{leancode}

\subsection{\texttt{R1CS.lean}: the relation and the size classes}

\lean{Sizes} with $\mathit{cols}=1+n_x+n_w$; the class shape \lean{classSizes k}
($m=\mathit{cols}=2^k$, $n_x=2^{10}$, $\mathit{nnz}\le2^{k+4}$) at \lean{classExp c}${}\in\{16,20,24\}$;
sparse matrices as entry lists (duplicates summed); the index with \lean{nnz\_le} on the list
lengths; $z=(1,x,w)$; \lean{R1CS} and \lean{rel}; and the harness padding maps, whose
satisfiability preservation in both directions is \lean{pad\_sat\_iff} in the library.
\emph{Check}: the constant wire is column 0 (\lean{Fin.cons 1}); \lean{SpMat.mulVec} sums all entries
of row $r$; \lean{R1CS} quantifies over every row; padding keeps public inputs at their columns,
moves witness columns past the larger $n_x$, and pads $x$ with zeros. A statement is a
representation: equal matrices given by different lists are different statements, so the harness
fixes a canonical form for benchmark indices.

\begin{leancode}[fontsize=\scriptsize,label={\footnotesize Contract/R1CS.lean}]
import RecursionFast.Contract.Field
import Mathlib.Data.Fin.Tuple.Basic

/-- Sizes of an R1CS instance: `m` rows, `nx` public inputs, `nw` witness entries, at most
`nnz` nonzero matrix entries in total. Columns: `cols = 1 + nx + nw` (the constant wire first). -/
structure Sizes where
  m : ℕ
  nx : ℕ
  nw : ℕ
  nnz : ℕ
  deriving DecidableEq

/-- Number of columns, `1 + nx + nw`. -/
abbrev Sizes.cols (sz : Sizes) : ℕ := sz.nx + sz.nw + 1

/-- The fixed class shape for exponent `k` (meaningful for `k ≥ 10`): `m = cols = 2^k`. -/
def classSizes (k : ℕ) : Sizes := ⟨2 ^ k, 2 ^ 10, 2 ^ k - 1 - 2 ^ 10, 2 ^ (k + 4)⟩

/-- The three scored classes, by exponent: `classExp 0 = 16`, `classExp 1 = 20`, `classExp 2 = 24`. -/
def classExp (c : Fin 3) : ℕ := 16 + 4 * c.val

/-- A **sparse** matrix: a list of entries `(row, column, value)`, indices in range by typing.
Entries at a repeated position are **summed**: the matrix denoted is the sum of its entries
(`Library/R1CS.lean`: `SpMat.toMat`, and `mulVec` equals Mathlib's dense `Matrix.mulVec`). -/
abbrev SpMat (sz : Sizes) := List (Fin sz.m × Fin sz.cols × K)

/-- Row `r` of the matrix-vector product, `Σ_{(r, c, v) ∈ M} v · z[c]` in K (left to right; the
order is irrelevant since `+` is xor). -/
def SpMat.mulVec {sz : Sizes} (M : SpMat sz) (z : Fin sz.cols → K) (r : Fin sz.m) : K :=
  M.foldl (fun acc e => if e.1 = r then acc + e.2.2 * z e.2.1 else acc) 0

/-- An R1CS index (the circuit): three sparse matrices with at most `nnz` entries in total (so an
indexer runs in time linear in `nnz`, not in `m · cols`). -/
structure Index (sz : Sizes) where
  A : SpMat sz
  B : SpMat sz
  C : SpMat sz
  nnz_le : A.length + B.length + C.length ≤ sz.nnz

/-- Public input `x ∈ K^nx`. -/
abbrev PubInput (sz : Sizes) := Fin sz.nx → K
/-- Witness `w ∈ K^nw`. -/
abbrev Wit (sz : Sizes) := Fin sz.nw → K

/-- The full assignment `z = (1, x, w)`. -/
def zVec {sz : Sizes} (x : PubInput sz) (w : Wit sz) : Fin sz.cols → K := Fin.cons 1 (Fin.append x w)

/-- **The relation**: every row satisfies `(A z)_r · (B z)_r = (C z)_r`. -/
def R1CS {sz : Sizes} (i : Index sz) (x : PubInput sz) (w : Wit sz) : Prop :=
  ∀ r, i.A.mulVec (zVec x w) r * i.B.mulVec (zVec x w) r = i.C.mulVec (zVec x w) r

instance {sz : Sizes} (i : Index sz) (x : PubInput sz) (w : Wit sz) : Decidable (R1CS i x w) :=
  inferInstanceAs (Decidable (∀ _, _))

/-- The relation as a set: statement `(index, public input)`, witness `w`. -/
def rel (sz : Sizes) : Set ((Index sz × PubInput sz) × Wit sz) := {p | R1CS p.1.1 p.1.2 p.2}

/-- Harness padding of an instance of size `sz` into a larger size `sz'` (column map): the
constant column stays at 0, public input `j < nx` stays at column `1 + j`, witness `j < nw` moves
to column `1 + nx' + j`; `none` if it does not fit. Rows keep their number. The padded public
input is `x` followed by zeros. (Satisfiability preservation and reflection: `Library/Pad.lean`.) -/
def padCol (sz sz' : Sizes) (c : Fin sz.cols) : Option (Fin sz'.cols) :=
  let t := if c.val ≤ sz.nx then c.val else c.val - sz.nx + sz'.nx
  if h : t < sz'.cols then some ⟨t, h⟩ else none

/-- The padded matrix: every entry moved by `padCol` (entries that do not fit are dropped). -/
def SpMat.pad {sz : Sizes} (sz' : Sizes) (M : SpMat sz) : SpMat sz' :=
  M.filterMap fun e => if h : e.1.val < sz'.m then (padCol sz sz' e.2.1).map (⟨e.1.val, h⟩, ·, e.2.2) else none

/-- The padded public input: `x` followed by zeros. -/
def PubInput.pad {sz : Sizes} (sz' : Sizes) (x : PubInput sz) : PubInput sz' :=
  fun j => if h : j.val < sz.nx then x ⟨j.val, h⟩ else 0
\end{leancode}

\subsection{\texttt{Shapes.lean}: round shapes, transcripts and oracles}

The fixed round menu (\cref{subsec:proto-rounds}): a prover round sends a raw vector of
$2^{\mathit{logLen}}$ elements of $K$, read in leaves of $2^{\mathit{logLeaf}}$; a verifier round
sends $n_E$ elements of $E$ and one $b$-bit index per $b$ in \lean{idx}. Transcript prefixes are
dependent functions on \lean{Fin m}, as in ArkLib. \emph{Check}: message types are exactly
\lean{OVal o} (no submitter type, so \enquote{for all transcripts} is \enquote{for all oracle
contents}); there are no $K$-valued challenges; \lean{Round.tapeWords} writes an $E$ element as its
three limbs and an index zero-extended to one word, in round order; prover messages become QUERY
oracles $1,2,\dots$ in round order; \lean{Oracle.get} returns $0$ beyond the length, which
\code{QUERY} never reads.

\begin{leancode}[fontsize=\scriptsize,label={\footnotesize Contract/Shapes.lean}]
import RecursionFast.Contract.Field
import Mathlib.Data.Fin.Tuple.Basic

/-- Contract round cap (DESIGN §3/§5): every protocol and reduction has at most `2^14` rounds. -/
def roundCap : ℕ := 2 ^ 14

/-- Shape of an oracle: `2^logLen` elements of K, read in leaves of `2^logLeaf` elements, so a
QUERY opens one of `2^(logLen - logLeaf)` leaves (Merkle depth `logLen - logLeaf`). -/
structure OShape where
  logLen : ℕ
  logLeaf : ℕ
  deriving DecidableEq

/-- Contents of an oracle of shape `o`. -/
abbrev OVal (o : OShape) := Fin (2 ^ o.logLen) → K

/-- A round of the contract's fixed menu. -/
inductive Round where
  /-- The prover sends an oracle of shape `o`. -/
  | msg (o : OShape)
  /-- The verifier sends `nE` elements of E and one `b`-bit index for each `b ∈ idx`. -/
  | chal (nE : ℕ) (idx : List ℕ)
  deriving DecidableEq

/-- What is sent in a round: a raw K-vector, or (E-vector, indices). -/
def Round.Ty : Round → Type
  | .msg o => OVal o
  | .chal nE idx => (Fin nE → E) × ((j : Fin idx.length) → BitVec (idx.get j))

/-- Whether the prover speaks in this round. -/
def Round.isMsg : Round → Bool
  | .msg _ => true
  | .chal _ _ => false

variable {n : ℕ} (sh : Fin n → Round)

/-- Transcript prefix of length `m`: the values of rounds `0, …, m-1`. -/
abbrev Transcript (m : Fin (n + 1)) := (i : Fin m) → (sh (Fin.castLE (Nat.le_of_lt_succ m.isLt) i)).Ty

/-- A full transcript: a value for every round. -/
abbrev FullTranscript := (i : Fin n) → (sh i).Ty

/-- The empty transcript. -/
def Transcript.empty : Transcript sh 0 := fun i => Fin.elim0 i

/-- Append the value of round `m` to a prefix of length `m`. -/
def Transcript.concat {sh : Fin n → Round} {m : Fin n} (tr : Transcript sh m.castSucc) (a : (sh m).Ty) :
    Transcript sh m.succ :=
  Fin.snoc (α := fun i : Fin (m.val + 1) => (sh (Fin.castLE m.isLt i)).Ty) tr a

/-- Indices of prover rounds. -/
abbrev MessageIdx := {m : Fin n // (sh m).isMsg = true}
/-- Indices of verifier (challenge) rounds. -/
abbrev ChallengeIdx := {m : Fin n // (sh m).isMsg = false}
/-- A choice of all challenges. -/
abbrev Challenges := (c : ChallengeIdx sh) → (sh c.1).Ty

/-- An oracle as the machine sees it: a shape and contents. -/
structure Oracle where
  shape : OShape
  data : OVal shape

instance : Inhabited Oracle := ⟨⟨⟨0, 0⟩, fun _ => 0⟩⟩

/-- Element `p` of an oracle (`0` beyond its length; QUERY never reads there). -/
def Oracle.get (o : Oracle) (p : ℕ) : K := if h : p < 2 ^ o.shape.logLen then o.data ⟨p, h⟩ else 0

/-- Number of leaves, `2^(logLen - logLeaf)`. -/
def OShape.leaves (o : OShape) : ℕ := 2 ^ (o.logLen - o.logLeaf)

/-- Leaf `l`: elements `l·2^logLeaf, …, (l+1)·2^logLeaf - 1`. -/
def Oracle.leaf (o : Oracle) (l : ℕ) : List K :=
  List.ofFn fun t : Fin (2 ^ o.shape.logLeaf) => o.get (l * 2 ^ o.shape.logLeaf + t.val)

/-- The oracles a round contributes (one for a `msg` round). -/
def Round.oracles : (r : Round) → r.Ty → List Oracle
  | .msg o, v => [⟨o, v⟩]
  | .chal _ _, _ => []

/-- The tape words a round contributes: E elements as 3 limbs each, then each index
zero-extended to one word (`b ≤ 64` is an admission rule). Prover rounds contribute none. -/
def Round.tapeWords : (r : Round) → r.Ty → List Word
  | .msg _, _ => []
  | .chal nE idx, (es, is) =>
      (List.ofFn fun j : Fin nE => (es j).toWords).flatten ++
        List.ofFn fun j : Fin idx.length => (is j).setWidth 64

/-- The prover messages of a transcript, as oracles `1, 2, …` in round order. -/
def FullTranscript.oracles {sh : Fin n → Round} (tr : FullTranscript sh) : List Oracle :=
  (List.ofFn fun i => (sh i).oracles (tr i)).flatten

/-- The challenges of a transcript, as tape words in round order. -/
def FullTranscript.tapeWords {sh : Fin n → Round} (tr : FullTranscript sh) : List Word :=
  (List.ofFn fun i => (sh i).tapeWords (tr i)).flatten
\end{leancode}

\subsection{\texttt{Security.lean}: completeness and round-by-round knowledge soundness}

The definitions of \cref{sec:security}, over the contract's round shapes and Mathlib only: a
deterministic \lean{Strategy}, a pure verifier, the honest transcript built by
\lean{Fin.induction}, \lean{PerfectComplete} for every challenge vector, the backward
\lean{Extractor}, the knowledge state function with its three laws, \lean{KSF.Bad}, \lean{RbrKS} in
counting form and \lean{SecureWith} with the error cap. \emph{Check}: \lean{Transcript sh m} is the
prefix of length $m$, so \lean{m.castSucc} and \lean{m.succ} are before and after round $m$;
\lean{cast E.eqIn} is an identity coercion; the $\exists w$ in \lean{Bad} ranges over next-stage
witnesses; the worst case is over every statement, challenge round and prefix; \lean{Nat.card} is
the size of a finite, non-empty challenge type; \lean{errorCap} is $2^{-128}$ (control
\lean{errorCap\_lt\_one}). At L1 \lean{RbrKS} is equivalent to RBR soundness
(\cref{thm:sec-collapse}).

\begin{leancode}[fontsize=\scriptsize,label={\footnotesize Contract/Security.lean}]
import RecursionFast.Contract.Shapes
import Mathlib.Basic.NNReal.Defs
import Mathlib.SetTheory.Cardinal.Finite

open scoped NNReal

/-- Contract error cap (DESIGN §5): every round's error is at most `2^-128`. -/
noncomputable def errorCap : ℝ≥0 := (2 : ℝ≥0)⁻¹ ^ 128

variable {n : ℕ} (sh : Fin n → Round) (Stmt Wit StmtOut WitOut : Type)

/-- A deterministic prover: each message is a function of the input and the transcript so far;
the output (statement, witness) is a function of the full transcript. -/
structure Strategy where
  next : Stmt → Wit → (m : MessageIdx sh) → Transcript sh m.1.castSucc → (sh m.1).Ty
  output : Stmt → Wit → FullTranscript sh → StmtOut × WitOut

/-- A pure verifier: a total function of the statement and the full transcript; `none` = reject. -/
abbrev PureVerifier := Stmt → FullTranscript sh → Option StmtOut

variable {sh Stmt Wit StmtOut WitOut}

/-- The honest transcript prefixes on challenge vector `ρ`: round `j` is the prover's message if
`j` is a prover round, and `ρ`'s entry otherwise. -/
def Strategy.prefixes (P : Strategy sh Stmt Wit StmtOut WitOut) (s : Stmt) (w : Wit)
    (ρ : Challenges sh) : (m : Fin (n + 1)) → Transcript sh m :=
  Fin.induction (Transcript.empty sh) fun j tr =>
    tr.concat (if h : (sh j).isMsg then P.next s w ⟨j, h⟩ tr else ρ ⟨j, by simpa using h⟩)

/-- The full honest transcript on challenge vector `ρ`. -/
def Strategy.transcript (P : Strategy sh Stmt Wit StmtOut WitOut) (s : Stmt) (w : Wit)
    (ρ : Challenges sh) : FullTranscript sh :=
  P.prefixes s w ρ (Fin.last n)

/-- **Perfect completeness**: for every `(s, w) ∈ relIn` and every challenge vector, the verifier
accepts the honest transcript with the prover's output statement, and the output pair is in
`relOut`. -/
def PerfectComplete (relIn : Set (Stmt × Wit)) (relOut : Set (StmtOut × WitOut))
    (P : Strategy sh Stmt Wit StmtOut WitOut) (V : PureVerifier sh Stmt StmtOut) : Prop :=
  ∀ s w, (s, w) ∈ relIn → ∀ ρ : Challenges sh,
    V s (P.transcript s w ρ) = some (P.output s w (P.transcript s w ρ)).1 ∧
      P.output s w (P.transcript s w ρ) ∈ relOut

variable (sh Stmt Wit WitOut) in
/-- A round-by-round extractor with intermediate witness types `WitMid` (ArkLib's
`Extractor.RoundByRound`): `extractMid m` maps a witness after round `m` to one before it. -/
structure Extractor (WitMid : Fin (n + 1) → Type) where
  eqIn : WitMid 0 = Wit
  extractMid : (m : Fin n) → Stmt → Transcript sh m.succ → WitMid m.succ → WitMid m.castSucc
  extractOut : Stmt → FullTranscript sh → WitOut → WitMid (Fin.last n)

/-- A **knowledge state function** for `V` and the named extractor `E`. Laws: `empty` (before
round 0 the state is the input relation), `next` (a prover message cannot make the state true
unless it was true before, with the extracted witness), `full` (an accepted full transcript with
an output pair in `relOut` makes the state true with the extracted final witness). -/
structure KSF (relIn : Set (Stmt × Wit)) (relOut : Set (StmtOut × WitOut))
    (V : PureVerifier sh Stmt StmtOut) {WitMid : Fin (n + 1) → Type}
    (E : Extractor sh Stmt Wit WitOut WitMid) where
  toFun : (m : Fin (n + 1)) → Stmt → Transcript sh m → WitMid m → Prop
  empty : ∀ s w, (s, cast E.eqIn w) ∈ relIn ↔ toFun 0 s (Transcript.empty sh) w
  next : ∀ (m : MessageIdx sh) s (tr : Transcript sh m.1.castSucc) (msg : (sh m.1).Ty) w,
    toFun m.1.succ s (tr.concat msg) w → toFun m.1.castSucc s tr (E.extractMid m.1 s (tr.concat msg) w)
  full : ∀ s tr sOut wOut, V s tr = some sOut → (sOut, wOut) ∈ relOut →
    toFun (Fin.last n) s tr (E.extractOut s tr wOut)

variable {relIn : Set (Stmt × Wit)} {relOut : Set (StmtOut × WitOut)}
  {V : PureVerifier sh Stmt StmtOut} {WitMid : Fin (n + 1) → Type}
  {E : Extractor sh Stmt Wit WitOut WitMid}

/-- A fresh challenge `ch` in round `c` is **bad** for prefix `tr` if some witness makes the state
true after `ch` while the extracted witness does not make it true before `ch`. -/
def KSF.Bad (kSF : KSF relIn relOut V E) (c : ChallengeIdx sh) (s : Stmt)
    (tr : Transcript sh c.1.castSucc) (ch : (sh c.1).Ty) : Prop :=
  ∃ w, ¬ kSF.toFun c.1.castSucc s tr (E.extractMid c.1 s (tr.concat ch) w) ∧
    kSF.toFun c.1.succ s (tr.concat ch) w

/-- **Worst-case RBR knowledge soundness with errors `ε`**, counting form: for every statement,
challenge round and prefix, at most an `ε c`-fraction of the round's challenges is bad. (All
challenge types are finite and nonempty, `Library/Field.lean`, so `Nat.card` is the size.) -/
def RbrKS (kSF : KSF relIn relOut V E) (ε : ChallengeIdx sh → ℝ≥0) : Prop :=
  ∀ s (c : ChallengeIdx sh) (tr : Transcript sh c.1.castSucc),
    (Nat.card {ch // kSF.Bad c s tr ch} : ℝ≥0) ≤ ε c * Nat.card (sh c.1).Ty

/-- The security clause shared by protocols and reductions: perfect completeness, an RBR-KS
state function for the named extractor at errors `ε`, and the error cap `ε c ≤ 2^-128`. -/
structure SecureWith (relIn : Set (Stmt × Wit)) (relOut : Set (StmtOut × WitOut))
    (P : Strategy sh Stmt Wit StmtOut WitOut) (V : PureVerifier sh Stmt StmtOut)
    (E : Extractor sh Stmt Wit WitOut WitMid) (ε : ChallengeIdx sh → ℝ≥0) : Prop where
  complete : PerfectComplete relIn relOut P V
  sound : ∃ kSF : KSF relIn relOut V E, RbrKS kSF ε
  cap : ∀ c, ε c ≤ errorCap
\end{leancode}

\subsection{\texttt{Protocol.lean}: holographic protocols}

A base-track protocol for one size: at most \lean{roundCap} rounds of the menu, a key of at most
\lean{vkCap}${}=2^{10}$ elements, an index oracle of a declared shape, an honest prover, the spec
verifier's decision from $(\mathsf{VK}\,i,x,\mathsf{IDX}\,i,\mathrm{tr})$, a named extractor and
declared errors. \emph{Check}: \lean{verify} never receives the index $i$; \lean{Secure} uses the
contract relation \lean{rel sz}, output relation \lean{Set.univ} and the declared errors.
Computability of the badge-path fields \lean{VK}, \lean{IDX}, \lean{prover} and \lean{extractor}
is a gate check (rule~4), not a Lean property; since round~4 the other fields, \lean{ε} included, may
be \lean{noncomputable} (A7-M2, \code{NOTICE-P2e.md}). The check only guarantees that compiled code
exists: a brute-force indexer is computable, and feasibility is the badge path's job (A7-M1). The
index-oracle cap $2^{k+8}$ is an admission clause in \code{Cost.lean}.

\begin{leancode}[fontsize=\scriptsize,label={\footnotesize Contract/Protocol.lean}]
import RecursionFast.Contract.Security
import RecursionFast.Contract.R1CS

open scoped NNReal

/-- Contract cap on the verification key: at most `2^10` K-elements. -/
def vkCap : ℕ := 2 ^ 10

/-- A holographic public-coin IOP for `rel sz`. -/
structure Protocol (sz : Sizes) where
  /-- Number of rounds, at most `roundCap = 2^14`. -/
  n : ℕ
  n_le : n ≤ roundCap
  /-- The round shapes (contract menu). -/
  shape : Fin n → Round
  /-- Verification key length and the (honest, computable) key. -/
  vkLen : ℕ
  vkLen_le : vkLen ≤ vkCap
  VK : Index sz → (Fin vkLen → K)
  /-- The index oracle (oracle 0): its shape and the (honest, computable) indexer. -/
  idxShape : OShape
  IDX : Index sz → OVal idxShape
  /-- The honest prover. -/
  prover : Strategy shape (Index sz × PubInput sz) (Wit sz) Unit Unit
  /-- The spec verifier's decision, from `VK i`, `x`, `IDX i` and the full transcript. -/
  verify : (Fin vkLen → K) → PubInput sz → OVal idxShape → FullTranscript shape → Bool
  /-- The named round-by-round extractor. -/
  WitMid : Fin (n + 1) → Type
  extractor : Extractor shape (Index sz × PubInput sz) (Wit sz) Unit WitMid
  /-- Declared per-round errors. -/
  ε : ChallengeIdx shape → ℝ≥0

namespace Protocol
variable {sz : Sizes} (P : Protocol sz)

/-- The verifier as a pure verifier on the statement `(i, x)`; `i` enters only via `VK i`, `IDX i`. -/
def specVerifier : PureVerifier P.shape (Index sz × PubInput sz) Unit :=
  fun s tr => if P.verify (P.VK s.1) s.2 (P.IDX s.1) tr then some () else none

/-- **Security of a protocol**: perfect completeness for `rel sz`, worst-case RBR knowledge
soundness with the named extractor at the declared errors, and `ε ≤ 2^-128` in every round
(the round cap `n ≤ 2^14` is a field of `Protocol`). -/
def Secure : Prop :=
  SecureWith (rel sz) Set.univ P.prover P.specVerifier P.extractor P.ε

end Protocol
\end{leancode}

\subsection{\texttt{Timer.lean}: the latency board's pipeline timer}

New in v0.3 (\cref{subsec:cost-lat}): a deterministic in-order dual-issue scoreboard. An operation
issues in the first free slot at or after the cycle in which its sources are ready, its result is
ready \lean{lat} cycles later, and a bubble \lean{pen} (taken branch or jump: 5; \code{QUERY}: 1,560
cycles per Merkle level plus the leaf length) flushes the rest of the cycle. \emph{Check}: the
timer is pure accounting and is never read by the decision; the constants match
\cref{subsec:cost-lat}. Its composition theorem $\mathit{time}(a\mathbin{+\!\!+}b)\le
\mathit{time}\,a+\mathit{time}\,b+L_{\max}$ is in \code{Library/Timer.lean}.

\begin{leancode}[fontsize=\scriptsize,label={\footnotesize Contract/Timer.lean}]
/-- One scheduled operation: destination register (none for `x0`/no result), source registers,
result latency in cycles, and the issue bubble after it in cycles. -/
structure Op where
  dst : Option Nat
  srcs : List Nat
  lat : Nat
  pen : Nat

/-- Scoreboard: the next free issue slot, and per register the cycle its value is ready. -/
structure SB where
  slot : Nat
  ready : Nat → Nat

/-- A fresh pipeline whose clock starts at cycle `t` (every register ready at `t`). -/
def SB.start (t : Nat) : SB := ⟨2 * t, fun _ => t⟩

/-- Elapsed cycles: the cycle after the last occupied slot. -/
def SB.time (sb : SB) : Nat := (sb.slot + 1) / 2

/-- Issue one operation (in order: never before `slot`; after its sources are ready). -/
def SB.step (sb : SB) (o : Op) : SB :=
  let t := o.srcs.foldl (fun m r => max m (2 * sb.ready r)) sb.slot
  { slot := if o.pen = 0 then t + 1 else 2 * (t / 2 + 1 + o.pen)
    ready := fun r => if o.dst = some r then t / 2 + o.lat else sb.ready r }

/-- Stall until cycle `t` (a `JOIN` waiting for its child). -/
def SB.wait (sb : SB) (t : Nat) : SB := { sb with slot := max sb.slot (2 * t) }

/-- Maximum result latency (the composition slack). -/
def Lmax : Nat := 16

/-- Bubble after a taken branch or a jump (D4 §2b `inord`: mispredict 5, static not-taken). -/
def TAKEN_PENALTY : Nat := 5

/-- Pipeline cycles per Merkle level of a QUERY on the latency board (D4 §1.2: software
BLAKE2s compression on the in-order model, 1,560 cycles). -/
def QUERY_LEVEL_CYCLES : Nat := 1560

/-- Fork-join overheads in cycles (D4 §3.2: a pooled runtime, about 1 µs at 2 GHz). -/
def C_SPAWN : Nat := 2000
def C_JOIN : Nat := 2000
\end{leancode}

\subsection{\texttt{Machine.lean}: the RISC-V oracle machine}

The largest file (297 code lines; v0.2: 165). Base instruction semantics and decoder are vendored;
the file adds the constants and memory map, the image, the 18-counter \lean{CostVec}, the decoder
arms for \code{ECALL}, Zbc and the W-ops, their semantics, the memory checks, \code{QUERY} in
closed form, the cost classes and pipeline operations, the shared accounting with first-touch
sets, fork-join harts, \lean{execute}, the initial state and \lean{run}. The module header
(omitted) is the specification restated in \cref{subsec:proto-machine}: system calls by \code{t0}
(0 \code{HALT}, 1 \code{QUERY}, 2 \code{HASH} reserved and trapping, 3 \code{SPAWN}, 4 \code{JOIN}
only with \lean{img.harts}), the memory map, fuel $2^{36}$ per hart with exhaustion rejecting, and
the separation of decision and accounting. \emph{Check}: \lean{decodeInstruction} accepts
\code{ECALL} only as the exact word and never \code{EBREAK} or CSR instructions; the Zbc
(\code{funct7}${}=5$) and W-op encodings; W-op division by zero and shift masking follow the ISA;
loads and stores are naturally aligned, inside $2^{36}$ bytes, and stores miss the tape window;
\code{QUERY} writes exactly the leaf and checks oracle, leaf and destination; \lean{HALT} accepts
only with $a_0=1$ and no live child; in \lean{execute} no \lean{if} or \lean{match} reads the
accounting \lean{x} or the timer \lean{sb}, except the spawn counter \lean{xc.spawns} (handles and
the $2^{16}$ cap); \lean{Ctx.allows} implements the hart discipline (own region, live children's
regions); a child that fails rejects the run; \lean{Acct.vec} counts a written page twice;
\lean{initState} places data at 0, the tape at \lean{TAPE\_BASE}, $\mathit{sp}=2^{36}$. Fuel is per
hart: a child runs with its parent's remaining fuel and the parent is not charged for the child's
steps, so a run with harts may execute up to $(\#\mathrm{spawns}+1)\cdot2^{36}$ steps (A6-F1,
row H12 of \cref{tab:thr-catalogue}).

\begin{leancode}[fontsize=\scriptsize,label={\footnotesize Contract/Machine.lean}]
import RiscvZkvm.Rv64.Execution
import RiscvZkvm.Interpreter.Decode
import RecursionFast.Contract.Shapes
import RecursionFast.Contract.Timer
import Mathlib.Data.Finset.Card
import Mathlib.Order.Interval.Finset.Nat

open RiscvZkvm.Rv64 (MachineState Instr Reg execInstrBr signExtend12)
open RiscvZkvm.Interpreter (decode regOfBits)

/-! ### Constants, image, counters -/

/-- Data memory size: addresses `[0, 2^36)` (bytes). -/
def MEM_BYTES : ℕ := 2 ^ 36
/-- Address of the first instruction (separate code space: loads and stores never reach it). -/
def CODE_BASE : ℕ := 0x1000
/-- Start of the read-only tape. -/
def TAPE_BASE : ℕ := 2 ^ 32
/-- Size of the read-only window at `TAPE_BASE` (= the tape admission cap, 1 MiB). -/
def TAPE_CAP : ℕ := 2 ^ 20
/-- Where a T2 program leaves its output accumulator's plain part (one word per K-element). -/
def OUT_BASE : ℕ := 2 ^ 33
/-- Step limit per hart; running out rejects. -/
def FUEL : ℕ := 2 ^ 36

/-- A program: raw instruction words, data words loaded at address 0, and whether it uses
fork-join harts (`SPAWN`/`JOIN`; if not, they trap as in v0.2). -/
structure Image where
  code : List (BitVec 32)
  data : List Word
  harts : Bool := false

/-- `k` little-endian `w`-bit words packed into the natural `x` (how the challenge renderer
encodes the submitted image files). -/
def unpack (w : ℕ) : ℕ → ℕ → List ℕ
  | 0, _ => []
  | k + 1, x => x % 2 ^ w :: unpack w k (x / 2 ^ w)

/-- Size in bytes (admission cap `2^20`, `Certificate.lean`). -/
def Image.byteSize (img : Image) : ℕ := 4 * img.code.length + 8 * img.data.length

/-- The 18 counters of a run (DESIGN §15.1, D4 §2a, §4.2): instructions by class (`add`:
add/sub/addi/lui/auipc and W forms; `logic`: and/or/xor/slt*; `shift`; `mul`: MUL/MULH*/MULW;
`div`: DIV*/REM* and W forms; `clmul`: Zbc; `ld`/`st`: doubleword, `ldsub`/`stsub`: narrower;
`br`: conditional branches; `jmp`: JAL/JALR), `sys` (HALT, SPAWN, JOIN), first-touched
doublewords `mWords` and 1 KiB pages `mPages` (code and data; a written page counts twice),
and QUERY count / Merkle levels / leaf bytes. -/
structure CostVec where
  add : ℕ := 0
  logic : ℕ := 0
  shift : ℕ := 0
  mul : ℕ := 0
  div : ℕ := 0
  clmul : ℕ := 0
  ld : ℕ := 0
  ldsub : ℕ := 0
  st : ℕ := 0
  stsub : ℕ := 0
  br : ℕ := 0
  jmp : ℕ := 0
  sys : ℕ := 0
  mWords : ℕ := 0
  mPages : ℕ := 0
  qCount : ℕ := 0
  qLevels : ℕ := 0
  qBytes : ℕ := 0
  deriving DecidableEq, Repr

/-- Instruction cost classes (the first twelve counters). -/
inductive Kind where
  | add | logic | shift | mul | div | clmul | ld | ldsub | st | stsub | br | jmp

/-- Count one instruction of class `k`. -/
def CostVec.tick (c : CostVec) : Kind → CostVec
  | .add => { c with add := c.add + 1 } | .logic => { c with logic := c.logic + 1 }
  | .shift => { c with shift := c.shift + 1 } | .mul => { c with mul := c.mul + 1 }
  | .div => { c with div := c.div + 1 } | .clmul => { c with clmul := c.clmul + 1 }
  | .ld => { c with ld := c.ld + 1 } | .ldsub => { c with ldsub := c.ldsub + 1 }
  | .st => { c with st := c.st + 1 } | .stsub => { c with stsub := c.stsub + 1 }
  | .br => { c with br := c.br + 1 } | .jmp => { c with jmp := c.jmp + 1 }

/-- Count one QUERY to an oracle of shape `o`: one Merkle path of depth `logLen - logLeaf` and
`8·2^logLeaf` leaf bytes. -/
def CostVec.addQuery (c : CostVec) (o : OShape) : CostVec :=
  { c with qCount := c.qCount + 1, qLevels := c.qLevels + (o.logLen - o.logLeaf),
           qBytes := c.qBytes + 8 * 2 ^ o.logLeaf }

/-! ### Decoding -/

/-- W-suffixed operations (upstream decodes only `ADDIW`). -/
inductive WordOp where
  | add | sub | sll | srl | sra | mul | div | divu | rem | remu
  deriving DecidableEq

/-- Zbc carry-less multiplications. -/
inductive ClOp where
  | clmul | clmulh | clmulr
  deriving DecidableEq

/-- A decoded instruction: upstream's, a W-op (register or immediate shift), or Zbc. -/
inductive Instruction where
  | base (i : Instr)
  | word (op : WordOp) (rd rs1 rs2 : Reg)
  | wordShift (op : WordOp) (rd rs1 : Reg) (shamt : BitVec 5)
  | cl (op : ClOp) (rd rs1 rs2 : Reg)
  deriving DecidableEq

/-- Decode a raw word: `ECALL` exactly (no `EBREAK`, no CSR), Zbc (`OP`, `funct7 = 5`,
`funct3 = 1/3/2`), the W-ops (`OP-32`, and `OP-IMM-32` shifts), then upstream `decode`. -/
def decodeInstruction (w : BitVec 32) : Option Instruction :=
  let opcode := (w.extractLsb' 0 7).toNat
  let rd := regOfBits (w.extractLsb' 7 5)
  let rs1 := regOfBits (w.extractLsb' 15 5)
  let rs2 := regOfBits (w.extractLsb' 20 5)
  let f3 := (w.extractLsb' 12 3).toNat
  let f7 := (w.extractLsb' 25 7).toNat
  if opcode = 0x73 then (if w = 0x00000073 then some (.base .ECALL) else none)
  else if opcode = 0x33 ∧ f7 = 5 then
    (match f3 with | 1 => some ClOp.clmul | 3 => some .clmulh | 2 => some .clmulr | _ => none).map
      fun op => .cl op rd rs1 rs2
  else if opcode = 0x3b then
    (match f7, f3 with
      | 0, 0 => some WordOp.add | 0x20, 0 => some .sub | 0, 1 => some .sll | 0, 5 => some .srl
      | 0x20, 5 => some .sra | 1, 0 => some .mul | 1, 4 => some .div | 1, 5 => some .divu
      | 1, 6 => some .rem | 1, 7 => some .remu | _, _ => none).map fun op => .word op rd rs1 rs2
  else if opcode = 0x1b ∧ f3 ≠ 0 then
    (match f7, f3 with
      | 0, 1 => some WordOp.sll | 0, 5 => some .srl | 0x20, 5 => some .sra | _, _ => none).map
      fun op => .wordShift op rd rs1 (w.extractLsb' 20 5)
  else (decode w).map .base

/-- 32-bit result of a W-op (RISC-V: shifts use 5 bits; `x/0 = -1`, `x % 0 = x`). -/
def wordResult (op : WordOp) (a b : BitVec 32) : BitVec 32 :=
  match op with
  | .add => a + b | .sub => a - b
  | .sll => a <<< (b.toNat % 32) | .srl => a >>> (b.toNat % 32) | .sra => a.sshiftRight (b.toNat % 32)
  | .mul => a * b
  | .div => if b = 0 then BitVec.allOnes 32 else a.sdiv b
  | .divu => if b = 0 then BitVec.allOnes 32 else a / b
  | .rem => if b = 0 then a else a.srem b
  | .remu => if b = 0 then a else a % b

/-- Zbc: low, high, and reversed (bits 126..63) halves of the 128-bit carry-less product. -/
def clResult (op : ClOp) (a b : Word) : Word :=
  match op with
  | .clmul => (K.clmul a b).setWidth 64
  | .clmulh => (K.clmul a b >>> 64).setWidth 64
  | .clmulr => (K.clmul a b >>> 63).setWidth 64

/-! ### One step -/

/-- `[a, a + len)` lies in data memory. -/
def inMemory (a len : ℕ) : Bool := a + len ≤ MEM_BYTES

/-- `[a, a + len)` lies in data memory and misses the read-only tape window. -/
def writable (a len : ℕ) : Bool :=
  inMemory a len && (a + len ≤ TAPE_BASE || TAPE_BASE + TAPE_CAP ≤ a)

/-- Address, width and direction (`true` = store) of each load and store. -/
def memAccess (s : MachineState) : Instr → Option (Word × ℕ × Bool)
  | .LD _ r o => some (s.getReg r + signExtend12 o, 8, false)
  | .LW _ r o | .LWU _ r o => some (s.getReg r + signExtend12 o, 4, false)
  | .LH _ r o | .LHU _ r o => some (s.getReg r + signExtend12 o, 2, false)
  | .LB _ r o | .LBU _ r o => some (s.getReg r + signExtend12 o, 1, false)
  | .SD r _ o => some (s.getReg r + signExtend12 o, 8, true)
  | .SW r _ o => some (s.getReg r + signExtend12 o, 4, true)
  | .SH r _ o => some (s.getReg r + signExtend12 o, 2, true)
  | .SB r _ o => some (s.getReg r + signExtend12 o, 1, true)
  | _ => none

/-- One non-`ECALL` instruction; `none` is a fault (misaligned or out-of-range access, store
into the tape, `EBREAK`/CSR/pseudo-instructions, which the decoder never produces anyway). -/
def ordinaryStep (s : MachineState) : Instruction → Option MachineState
  | .base .ECALL | .base .EBREAK | .base (.CSRS _ _) | .base (.LI _ _) | .base (.MV _ _)
  | .base .NOP => none
  | .base i =>
    match memAccess s i with
    | some (a, width, store) =>
      if a.toNat % width = 0 && (if store then writable a.toNat width else inMemory a.toNat width)
      then some (execInstrBr s i) else none
    | none => some (execInstrBr s i)
  | .word op rd rs1 rs2 =>
    let v := wordResult op ((s.getReg rs1).truncate 32) ((s.getReg rs2).truncate 32)
    some ((s.setReg rd (v.signExtend 64)).setPC (s.pc + 4))
  | .wordShift op rd rs1 sh =>
    let v := wordResult op ((s.getReg rs1).truncate 32) (sh.zeroExtend 32)
    some ((s.setReg rd (v.signExtend 64)).setPC (s.pc + 4))
  | .cl op rd rs1 rs2 => some ((s.setReg rd (clResult op (s.getReg rs1) (s.getReg rs2))).setPC (s.pc + 4))

/-- Fetch at `pc` (4-aligned, at or above `CODE_BASE`) and decode. -/
def fetch (img : Image) (s : MachineState) : Option Instruction :=
  if CODE_BASE ≤ s.pc.toNat ∧ s.pc.toNat % 4 = 0 then
    (img.code[(s.pc.toNat - CODE_BASE) / 4]?).bind decodeInstruction
  else none

/-- QUERY (see the header). Memory after the call, in closed form: the destination words hold
the leaf, every other word is unchanged. Returns the new state and the oracle's shape. -/
def query (os : List Oracle) (s : MachineState) : Option (MachineState × OShape) := do
  let o ← os[(s.getReg .x10).toNat]?
  let l := (s.getReg .x11).toNat
  let dst := (s.getReg .x12).toNat
  let bytes := 8 * 2 ^ o.shape.logLeaf
  if l < o.shape.leaves ∧ dst % 8 = 0 ∧ writable dst bytes = true then
    some ({ s with
      mem := fun a => if dst ≤ a.toNat ∧ a.toNat < dst + bytes
        then (o.get (l * 2 ^ o.shape.logLeaf + (a.toNat - dst) / 8)).bits else s.mem a
      pc := s.pc + 4 }, o.shape)
  else none

/-! ### Cost classes and pipeline operations -/

/-- The raw word at `pc` (only consulted after `fetch` succeeded). -/
def rawAt (img : Image) (s : MachineState) : BitVec 32 :=
  (img.code[(s.pc.toNat - CODE_BASE) / 4]?).getD 0

/-- Cost class of an instruction (`ECALL` is counted as `sys`; `FENCE` as `add`). -/
def kind : Instruction → Kind
  | .base i =>
    match i with
    | .ADD .. | .SUB .. | .ADDI .. | .LUI .. | .AUIPC .. | .ADDIW .. => .add
    | .AND .. | .OR .. | .XOR .. | .SLT .. | .SLTU .. | .ANDI .. | .ORI .. | .XORI .. | .SLTI .. | .SLTIU .. =>
        .logic
    | .SLL .. | .SRL .. | .SRA .. | .SLLI .. | .SRLI .. | .SRAI .. => .shift
    | .MUL .. | .MULH .. | .MULHSU .. | .MULHU .. => .mul
    | .DIV .. | .DIVU .. | .REM .. | .REMU .. => .div
    | .LD .. => .ld
    | .LW .. | .LWU .. | .LH .. | .LHU .. | .LB .. | .LBU .. => .ldsub
    | .SD .. => .st
    | .SW .. | .SH .. | .SB .. => .stsub
    | .BEQ .. | .BNE .. | .BLT .. | .BGE .. | .BLTU .. | .BGEU .. => .br
    | .JAL .. | .JALR .. => .jmp
    | _ => .add
  | .word op .. => match op with | .add | .sub => .add | .sll | .srl | .sra => .shift | .mul => .mul | _ => .div
  | .wordShift .. => .shift
  | .cl .. => .clmul

/-- Result latency per class in cycles (D4 §2b `inord`, SiFive-7-like: loads, `mul`, `clmul` 3,
`div` 16, everything else 1). -/
def latency : Kind → ℕ
  | .mul | .clmul | .ld | .ldsub => 3
  | .div => 16
  | _ => 1

/-- The pipeline operation of a raw word: destination and sources from its format (`x0` is never
a destination), result latency `lat`, bubble `pen`. -/
def opOf (w : BitVec 32) (lat pen : ℕ) : Op :=
  let f (lo : ℕ) := (w.extractLsb' lo 5).toNat
  let d := if f 7 = 0 then none else some (f 7)
  let (dst, srcs) : Option ℕ × List ℕ := match (w.extractLsb' 0 7).toNat with
    | 0x33 | 0x3b => (d, [f 15, f 20])
    | 0x13 | 0x1b | 0x03 | 0x67 => (d, [f 15])
    | 0x23 | 0x63 => (none, [f 15, f 20])
    | 0x37 | 0x17 | 0x6f => (d, [])
    | _ => (none, [5, 10, 11, 12, 13, 14, 15, 16])
  ⟨dst, srcs, lat, pen⟩

/-- A system call reads `t0, a0, …, a6`; `SPAWN` writes `a0`. -/
def ecallOp (dst : Option ℕ) (pen : ℕ) : Op := ⟨dst, [5, 10, 11, 12, 13, 14, 15, 16], 1, pen⟩

/-! ### Accounting (shared by all harts of a run; never consulted by the decision) -/

/-- Key of a doubleword or page number `n`: data `2n`, code `2n+1` (Harvard: separate spaces). -/
def key (code : Bool) (n : ℕ) : ℕ := 2 * n + if code then 1 else 0

/-- Counters, first-touch sets (doublewords, pages, written pages), pipeline work `W`, spawns. -/
structure Acct where
  cost : CostVec := {}
  words : Finset ℕ := ∅
  pages : Finset ℕ := ∅
  dirty : Finset ℕ := ∅
  work : ℕ := 0
  spawns : ℕ := 0

/-- Touch the bytes `[a, a + len)` of code or data; a store also dirties their pages. -/
def Acct.touch (x : Acct) (code store : Bool) (a len : ℕ) : Acct :=
  let ps := (Finset.Ico (a / 1024) ((a + len + 1023) / 1024)).image (key code)
  { x with words := x.words ∪ (Finset.Ico (a / 8) ((a + len + 7) / 8)).image (key code)
           pages := x.pages ∪ ps, dirty := if store then x.dirty ∪ ps else x.dirty }

/-- One more `HALT`/`SPAWN`/`JOIN`. -/
def Acct.sys (x : Acct) : Acct := { x with cost := { x.cost with sys := x.cost.sys + 1 } }

/-- The 18-counter vector: first-touch counts are the sizes of the sets. -/
def Acct.vec (x : Acct) : CostVec :=
  { x.cost with mWords := x.words.card, mPages := x.pages.card + x.dirty.card }

/-- Issue `o` on the hart's pipeline and charge the cycles it adds as work. -/
def tock (sb : SB) (x : Acct) (o : Op) : SB × Acct :=
  let sb' := sb.step o
  (sb', { x with work := x.work + (sb'.time - sb.time) })

/-! ### Harts (D4 §3.2): fork-join with declared disjoint regions, sequentialised

`SPAWN` (`a0` = entry pc, `a1` = argument, `a2` = child `sp`, `a3`/`a4` = base/length of the
child's write region `Wr`, `a5`/`a6` = base/length of its read region `R`; all 8-aligned, `Wr`
writable) runs the child **to completion at once** (registers `0` except `a0 = a1`, `sp = a2`;
its memory view is the parent's; it may load only from `R ∪ Wr`, store and QUERY only into
`Wr`; it must end with `HALT(1)`, else the whole run rejects) and returns a handle in `a0`. The
child's writes become visible to the parent only at `JOIN` (`a0` = handle). While a child is
live, no other hart may store into its `R ∪ Wr` or load from its `Wr`; the regions of a new
child must be accessible to its parent under these rules. Every violation faults (= reject).
Hence every interleaving of live harts yields the same state at `JOIN`, so a deployed runtime
may run children in parallel. Caps: depth `8`, `64` live children per hart, `2^16` spawns. -/

/-- A byte interval `[base, base + len)`. -/
structure Region where
  base : ℕ
  len : ℕ

/-- `[a, a + len)` lies inside the region. -/
def Region.has (r : Region) (a len : ℕ) : Bool := r.base ≤ a && a + len ≤ r.base + r.len
/-- `[a, a + len)` is disjoint from the region. -/
def Region.misses (r : Region) (a len : ℕ) : Bool := a + len ≤ r.base || r.base + r.len ≤ a

/-- A spawned, not yet joined child: its regions, final memory and the cycle it finished. -/
structure Child where
  handle : ℕ
  wr : Region
  rd : Region
  mem : Word → Word
  finish : ℕ

/-- Hart context: own `(Wr, R)` (`none` = the root hart) and live children. -/
structure Ctx where
  own : Option (Region × Region) := none
  live : List Child := []
  depth : ℕ := 0

/-- May this hart load (`store = false`) or store `[a, a + len)`? -/
def Ctx.allows (h : Ctx) (store : Bool) (a len : ℕ) : Bool :=
  (match h.own with | none => true | some (wr, rd) => wr.has a len || (!store && rd.has a len)) &&
    h.live.all fun ch => ch.wr.misses a len && (!store || ch.rd.misses a len)

/-- Address, width and direction (`true` = store) of a load or store instruction. -/
def access (s : MachineState) : Instruction → Option (ℕ × ℕ × Bool)
  | .base i => (memAccess s i).map fun (a, w, st) => (a.toNat, w, st)
  | _ => none

/-! ### Execution -/

/-- Run one hart for at most `fuel` steps in context `h`, threading the pipeline `sb` and the
accounting `x`. `some` iff the hart reaches `HALT` with `a0 = 1` and no live child; every
fault, trap, rejecting `HALT`, failed child and fuel exhaustion is `none` (reject). -/
def execute (img : Image) (os : List Oracle) :
    ℕ → Ctx → MachineState → SB → Acct → Option (Acct × SB × MachineState)
  | 0, _, _, _, _ => none
  | fuel + 1, h, s, sb, x =>
    let x := x.touch true false s.pc.toNat 4
    match fetch img s with
    | none => none
    | some (.base .ECALL) =>
      match (s.getReg .x5).toNat with
      | 0 => if s.getReg .x10 = 1 ∧ h.live.isEmpty then
          let (sb', x') := tock sb x (ecallOp none 0); some (x'.sys, sb', s) else none
      | 1 => match query os s with
        | some (s', o) =>
          let dst := (s.getReg .x12).toNat
          if h.allows true dst (8 * 2 ^ o.logLeaf) then
            let (sb', x') := tock sb x (ecallOp none
              (QUERY_LEVEL_CYCLES * (o.logLen - o.logLeaf) + 2 ^ o.logLeaf))
            execute img os fuel h s' sb'
              ({ x' with cost := x'.cost.addQuery o }.touch false true dst (8 * 2 ^ o.logLeaf))
          else none
        | none => none
      | 3 =>
        let wr : Region := ⟨(s.getReg .x13).toNat, (s.getReg .x14).toNat⟩
        let rd : Region := ⟨(s.getReg .x15).toNat, (s.getReg .x16).toNat⟩
        let (sb', x') := tock sb x (ecallOp (some 10) 0)
        if img.harts ∧ h.depth < 8 ∧ h.live.length < 64 ∧ wr.base % 8 = 0 ∧ wr.len % 8 = 0 ∧
            rd.base % 8 = 0 ∧ rd.len % 8 = 0 ∧ writable wr.base wr.len ∧ inMemory rd.base rd.len ∧
            h.allows true wr.base wr.len ∧ h.allows false rd.base rd.len then
          let child : MachineState := { s with pc := s.getReg .x10, regs := fun r =>
            match r with | .x10 => s.getReg .x11 | .x2 => s.getReg .x12 | _ => 0 }
          match execute img os fuel ⟨some (wr, rd), [], h.depth + 1⟩ child
              (SB.start (sb'.time + C_SPAWN)) x'.sys with
          | some (xc, sbc, sc) =>
            if xc.spawns < 2 ^ 16 then
              execute img os fuel { h with live := ⟨xc.spawns, wr, rd, sc.mem, sbc.time⟩ :: h.live }
                ((s.setReg .x10 xc.spawns).setPC (s.pc + 4)) sb'
                { xc with spawns := xc.spawns + 1, work := xc.work + C_SPAWN + C_JOIN }
            else none
          | none => none
        else none
      | 4 =>
        match h.live.find? (·.handle = (s.getReg .x10).toNat) with
        | some ch =>
          let (sb', x') := tock sb x (ecallOp none 0)
          execute img os fuel { h with live := h.live.filter (·.handle ≠ ch.handle) }
            { s with pc := s.pc + 4, mem := fun a => if ch.wr.has a.toNat 8 then ch.mem a else s.mem a }
            (sb'.wait (ch.finish + C_JOIN)) x'.sys
        | none => none
      | _ => none
    | some i =>
      let acc := access s i
      if acc.all fun (a, w, st) => h.allows st a w then
        match ordinaryStep s i with
        | some s' =>
          let (sb', x') := tock sb x
            (opOf (rawAt img s) (latency (kind i)) (if s'.pc = s.pc + 4 then 0 else TAKEN_PENALTY))
          let x' := { x' with cost := x'.cost.tick (kind i) }
          execute img os fuel h s' sb' (match acc with | some (a, w, st) => x'.touch false st a w | none => x')
        | none => none
      else none

/-- Initial state: data words from address 0, the tape from `TAPE_BASE` (both zero beyond
their length), `sp = 2^36`, all other registers 0, `pc = CODE_BASE`. -/
def initState (img : Image) (tape : List Word) : MachineState where
  regs := fun r => match r with | .x2 => BitVec.ofNat 64 MEM_BYTES | _ => 0
  mem := fun a => if TAPE_BASE ≤ a.toNat then tape.getD ((a.toNat - TAPE_BASE) / 8) 0
    else img.data.getD (a.toNat / 8) 0
  pc := BitVec.ofNat 64 CODE_BASE

/-- What an accepting run yields: the 18 counters (summed over all harts: the zkVM work), the
pipeline work `W` (all harts, plus `C_SPAWN + C_JOIN` per spawn) and the span `S` (the root
hart's clock, which waits at each `JOIN` for its child). -/
structure Final where
  cost : CostVec
  work : ℕ
  span : ℕ

/-- **The deployed verifier with its accounting**: run `img` on `tape` with oracles `os`
(QUERY oracle `j` is `os[j]`). `some (final, state)` = accept. -/
def runFull (img : Image) (tape : List Word) (os : List Oracle) : Option (Final × MachineState) :=
  (execute img os FUEL {} (initState img tape) (SB.start 0) {}).map fun (x, sb, s) =>
    (⟨x.vec, x.work, sb.time⟩, s)

/-- **The deployed verifier** (decision, counters, final state). -/
def run (img : Image) (tape : List Word) (os : List Oracle) : Option (CostVec × MachineState) :=
  (runFull img tape os).map fun r => (r.1.cost, r.2)

/-- The `len` words a T2 program left at `OUT_BASE`. -/
def readOut (s : MachineState) (len : ℕ) : List Word :=
  List.ofFn fun j : Fin len => s.getMem (BitVec.ofNat 64 (OUT_BASE + 8 * j.val))
\end{leancode}

\subsection{\texttt{Agrees.lean}: the verifier--spec link}

The tape (\lean{sizes} $\|$ \lean{VK i} $\|$ $x$ $\|$ challenges), the oracle list (\lean{IDX i}
then the messages), \lean{Agrees} and the machine as a pure verifier for the transfer lemma.
\emph{Check}: \lean{Agrees} quantifies over every index of the size, every public input and every
full transcript, i.e.\ every raw oracle content and every challenge; preprocessing is the honest
\lean{VK i}, \lean{IDX i}; acceptance is \lean{isSome}, so faults, traps and fuel exhaustion are
rejections. No cost is transferred from the spec.

\begin{leancode}[fontsize=\scriptsize,label={\footnotesize Contract/Agrees.lean}]
import RecursionFast.Contract.Protocol
import RecursionFast.Contract.Machine

/-- Tape header: the four sizes `m, nx, nw, nnz`, one word each. -/
def Sizes.header (sz : Sizes) : List Word := [sz.m, sz.nx, sz.nw, sz.nnz].map (BitVec.ofNat 64)

/-- A K-vector as tape words, in index order. -/
def wordsOf {l : ℕ} (v : Fin l → K) : List Word := List.ofFn fun j => (v j).bits

namespace Protocol
variable {sz : Sizes} (P : Protocol sz)

/-- The tape: `sizes ‖ VK i ‖ x ‖ challenges` (challenges in round order, `Round.tapeWords`). -/
def tape (i : Index sz) (x : PubInput sz) (tr : FullTranscript P.shape) : List Word :=
  sz.header ++ wordsOf (P.VK i) ++ wordsOf x ++ tr.tapeWords

/-- The QUERY oracles: `IDX i` is oracle 0, the prover messages are oracles `1, 2, …`. -/
def oracles (i : Index sz) (tr : FullTranscript P.shape) : List Oracle :=
  ⟨P.idxShape, P.IDX i⟩ :: tr.oracles

/-- **Decision agreement** of the image with the spec verifier, on all transcripts. -/
def Agrees (img : Image) : Prop :=
  ∀ (i : Index sz) (x : PubInput sz) (tr : FullTranscript P.shape),
    (run img (P.tape i x tr) (P.oracles i tr)).isSome = P.verify (P.VK i) x (P.IDX i) tr

/-- The deployed machine as a pure verifier on `(i, x)` (for the transfer lemma,
`Library/Transfer.lean`: `Agrees` + `Secure` ⇒ the machine verifier is secure). -/
def machineVerifier (img : Image) : PureVerifier P.shape (Index sz × PubInput sz) Unit :=
  fun s tr => (run img (P.tape s.1 s.2 tr) (P.oracles s.1 tr)).map fun _ => ()

end Protocol
\end{leancode}

\subsection{\texttt{Cost.lean}: cost clauses, profiles, score functions and admission}\label{subsec:app-cost}

The componentwise order on the 18 counters; \lean{Accepting\allowbreak CostAtMost} and the optional
\lean{Accepting\allowbreak LatencyAtMost}; the four zkVM profiles and the displayed v0.2 weights; the per-run
cost $C_P$; the shape \lean{Summary} and its Fiat--Shamir, byte and blow-up terms;
\lean{scoreFourth}, \lean{scoreZk4} and \lean{scoreLat3}; and the admission predicates.
\emph{Check}: \lean{CostVec.toList} lists each of the 18 fields once, so \lean{≤} is componentwise;
the cost clause quantifies over honest preprocessing and all $x$ and transcripts, and only over
accepting runs; the 15 weights of each profile match \cref{tab:cost-profiles} (\lean{zip} stops at
the shorter list, so the three query counters are priced only through
$(P_H+35)\,q^{\mathrm{lev}}+P_H\,(q^{\mathrm{bytes}}/64+q^{\mathrm{count}})$, the leaf term an upper
bound with at most one compression of slack per \code{QUERY} (A6-P1), and a profile with fewer than
15 weights would silently drop counters, A6-N2); \lean{fs} charges absorbing
$\mathsf{VK}\,\|\,x$ (A4-F6); \lean{scoreZk4} is the product over 4 profiles and 3 classes of $S^4$,
so $\mathrm{Score}_{\mathrm{zk}}=\mathit{zk4}^{1/48}$; message oracles are capped at
$2^{10}\max(m,\mathit{cols})$ and the index oracle at $2^{8}\max(m,\mathit{cols})$
(\lean{IdxAdmissible}, A5-M4).

\begin{leancode}[fontsize=\scriptsize,label={\footnotesize Contract/Cost.lean}]
import RecursionFast.Contract.Agrees
import Mathlib.Algebra.Order.Floor.Defs
import Mathlib.Data.Rat.Floor
import Mathlib.Algebra.BigOperators.Fin

/-- The counters in a fixed order: `add, logic, shift, mul, div, clmul, ld, ldsub, st, stsub,
br, jmp, sys, mWords, mPages, qCount, qLevels, qBytes`. -/
def CostVec.toList (c : CostVec) : List ℕ :=
  [c.add, c.logic, c.shift, c.mul, c.div, c.clmul, c.ld, c.ldsub, c.st, c.stsub, c.br, c.jmp,
    c.sys, c.mWords, c.mPages, c.qCount, c.qLevels, c.qBytes]

/-- Componentwise order on counters. -/
instance : LE CostVec := ⟨fun a b => List.Forall₂ (· ≤ ·) a.toList b.toList⟩

instance (a b : CostVec) : Decidable (a ≤ b) := inferInstanceAs (Decidable (List.Forall₂ _ _ _))

/-- **Accepting-run cost bound**: every accepting run, for every index of the size, every `x`
and every transcript, has counters `≤ c` componentwise. -/
def Protocol.AcceptingCostAtMost {sz : Sizes} (P : Protocol sz) (img : Image) (c : CostVec) : Prop :=
  ∀ (i : Index sz) x (tr : FullTranscript P.shape) v s, run img (P.tape i x tr) (P.oracles i tr) = some (v, s) →
      v ≤ c

/-- **Accepting-run latency bound** (latency board, optional): pipeline work `≤ W` and span `≤ S`. -/
def Protocol.AcceptingLatencyAtMost {sz : Sizes} (P : Protocol sz) (img : Image) (W S : ℕ) : Prop :=
  ∀ (i : Index sz) x (tr : FullTranscript P.shape) f s,
    runFull img (P.tape i x tr) (P.oracles i tr) = some (f, s) → f.work ≤ W ∧ f.span ≤ S

/-! ### zkVM profiles (D4 §2a, normalised table; refreshed per season, re-scoring without re-proof) -/

/-- A zkVM cost profile: weights of `add … mPages` (the first 15 counters, `toList` order),
the hash price `PH` per 64-byte compression (a Merkle level, a 64-byte block of a QUERY leaf, a
transcript absorption), and the price `lam` per proof byte. -/
structure Profile where
  w : List ℚ
  PH : ℚ
  lam : ℚ

/-- ZK-A ZisK (`zisk_ops_costs.rs` areas, logic op = 1). `sys` = 1 in every profile ([J]: D4's
table has no `sys` row; v0.2 charged 1). -/
def zisk : Profile :=
  ⟨[73/100, 1, 97/100, 129/100, 129/100, 97/100, 6/10, 6/10, 6/10, 6/10, 1, 73/100, 1, 0, 0], 72, 1/8⟩
/-- ZK-B SP1 Hypercube (`rv64im_costs.json` areas, Add = 1; a hypothetical clmul chip priced as Mul). -/
def sp1 : Profile :=
  ⟨[1, 110/100, 127/100, 138/100, 264/100, 138/100, 105/100, 110/100, 105/100, 113/100, 109/100, 1, 1,
    12/100, 0], 174, 1/8⟩
/-- ZK-C RISC Zero (1 per instruction, bitwise and DIV/REM 2, 1,130 cycles per page in or out). -/
def risc0 : Profile := ⟨[1, 2, 1, 1, 2, 1, 1, 1, 1, 1, 1, 1, 1, 0, 1130], 100, 1/8⟩
/-- ZK-D Jolt (lookup-uniform; DIV/REM virtual sequences 8, sub-word load 3, sub-word store 9). -/
def jolt : Profile := ⟨[1, 1, 1, 1, 8, 1, 1, 3, 1, 9, 1, 1, 1, 0, 0], 200, 1/8⟩
/-- v0.2 / sig.golf weights (1; MUL = DIV = CLMUL = 4; `P_H = 8`): **displayed only**. -/
def uniformV02 : Profile := ⟨[1, 1, 1, 4, 4, 4, 1, 1, 1, 1, 1, 1, 1, 0, 0], 8, 1/8⟩

/-- The four scored profiles: ZisK, SP1 Hypercube, RISC Zero, Jolt. -/
def profiles : Fin 4 → Profile := ![zisk, sp1, risc0, jolt]

/-- `C_P = w_P · c + (P_H + 35)·q_levels + P_H·(q_bytes/64 + q_count)`: a QUERY is a virtual
Merkle opening, priced as its path (`P_H + 35` per level) plus the hash of the leaf it reads.
**Leaf hash** (A6-P1; v0.2 charged `q_bytes/8`, i.e. `P_H = 8`, in every profile): QUERY `i` reads
`b_i = 8·2^logLeaf` bytes and hashing them costs `⌈b_i/64⌉` compressions of 64 bytes; since
`⌈b/64⌉ ≤ b/64 + 1`, the sum over a run is `Σ_i ⌈b_i/64⌉ ≤ q_bytes/64 + q_count`. This upper bound
uses only the existing counters (exact up to one compression per QUERY; for leaves of `< 8`
elements it charges at most `P_H` per QUERY, where `⌈b/64⌉ = 1`), at the profile's own `P_H`, so a
leaf costs what absorbing the same words into the transcript costs (`Summary.fs`, A4-F6). -/
def Profile.C (P : Profile) (c : CostVec) : ℚ :=
  ((P.w.zip c.toList).map fun (a, n) => a * n).sum + (P.PH + 35) * c.qLevels +
    P.PH * ((c.qBytes : ℚ) / 64 + c.qCount)

/-! ### Declared shape: the score inputs a protocol's round shapes determine -/

variable {n : ℕ} (sh : Fin n → Round)

/-- Per-round sums over a shape. -/
def shapeSum (f : Round → ℕ) : ℕ := ((List.ofFn sh).map f).sum

/-- The shape-derived score inputs: rounds, prover rounds, E-challenge elements, index
challenges, total message length `Σ 2^logLen`, and tape input words (header ‖ VK ‖ x for a
protocol; header ‖ plain parts for an accumulation step). Rendered into the challenge. -/
structure Summary where
  n : ℕ
  msgs : ℕ
  nE : ℕ
  nIdx : ℕ
  msgLen : ℕ
  inWords : ℕ
  deriving DecidableEq

/-- The summary of a shape with `inWords` tape input words. -/
def summarize (inWords : ℕ) : Summary :=
  ⟨n, shapeSum sh (fun | .msg _ => 1 | _ => 0), shapeSum sh (fun | .chal nE _ => nE | _ => 0),
    shapeSum sh (fun | .chal _ idx => idx.length | _ => 0), shapeSum sh (fun | .msg o => 2 ^ o.logLen | _ => 0),
    inWords⟩

namespace Summary
variable (d : Summary)

/-- `FS_P = rounds·(2·P_H + 30) + (¾·P_H + 6)·#E + 4·#idx + ⌈(input words)·8/64⌉·P_H`: squeezing
per round and challenge, plus absorbing the input `VK ‖ x` (A4-F6); v0.2's 46/12/4 at `P_H = 8`. -/
def fs (P : Profile) : ℚ :=
  d.n * (2 * P.PH + 30) + (3 / 4 * P.PH + 6) * d.nE + 4 * d.nIdx + ((d.inWords - 4 + 7) / 8 : ℕ) * P.PH

/-- `B = tapeBytes + q_bytes + 32·q_levels + 32·#msg rounds` (proof-byte charge). -/
def bytes (c : CostVec) : ℕ := 8 * (d.inWords + 3 * d.nE + d.nIdx) + c.qBytes + 32 * c.qLevels + 32 * d.msgs

/-- `V = Σ_{msg rounds} ℓ_j / (N + nnz)` (declared message lengths; `IDX` excluded). -/
def blowup (sz : Sizes) : ℚ := d.msgLen / (max sz.m sz.cols + sz.nnz : ℕ)

/-- `S_P⁴ = (C_P + FS_P + ⌈λ_P·B⌉)⁴ · max(1, V/16)` (D4 §4.2), root-free. -/
def scoreFourth (P : Profile) (sz : Sizes) (c : CostVec) : ℚ :=
  (P.C c + d.fs P + ⌈P.lam * d.bytes c⌉) ^ 4 * max 1 (d.blowup sz / 16)

end Summary

/-- Size of class `c` (`k = 16, 20, 24`). -/
abbrev classSz (c : Fin 3) : Sizes := classSizes (classExp c)

/-- **Recursion board (T0/T1)**: `scoreZk4 = Π_P Π_k S_P,k⁴`, so
`Score_zk = (Π_P Π_k S_P,k)^(1/12) = scoreZk4^(1/48)` (lower is better). -/
def scoreZk4 (d : Fin 3 → Summary) (c : Fin 3 → CostVec) : ℚ :=
  ∏ P : Fin 4, ∏ k : Fin 3, (d k).scoreFourth (profiles P) (classSz k) (c k)

/-- **Latency board**: with `T_16,k = W_k/16 + S_k` (Brent, 16 cores),
`scoreLat3 = Π_k T_16,k` and `Score_lat = scoreLat3^(1/3)`. `lat k = (W_k, S_k)`. -/
def scoreLat3 (lat : Fin 3 → ℕ × ℕ) : ℚ := ∏ k : Fin 3, (((lat k).1 : ℚ) / 16 + (lat k).2)

/-! ### Admission caps (certificate clauses) -/

/-- An oracle shape is admissible at size `sz`: `leaf ≤ ℓ ≤ 2^10·N` (`= 2^(k+10)` at class `k`). -/
def OShape.Admissible (sz : Sizes) (o : OShape) : Prop :=
  o.logLeaf ≤ o.logLen ∧ 2 ^ o.logLen ≤ 2 ^ 10 * max sz.m sz.cols

/-- The index oracle is admissible at size `sz`: `leaf ≤ idxLen ≤ 2^8·N` (`= 2^(k+8)` at class `k`;
DESIGN §15.3, A5-M4: with it the compiled bound's `t*` is a function of declared quantities). -/
def OShape.IdxAdmissible (sz : Sizes) (o : OShape) : Prop :=
  o.logLeaf ≤ o.logLen ∧ 2 ^ o.logLen ≤ 2 ^ 8 * max sz.m sz.cols

/-- A round is admissible: admissible oracle shape, index challenges of at most 64 bits. -/
def Round.Admissible (sz : Sizes) : Round → Prop
  | .msg o => o.Admissible sz
  | .chal _ idx => ∀ b ∈ idx, b ≤ 64

/-- Shape caps shared by every protocol and reduction: rounds admissible, tape `≤ 2^20` bytes,
blow-up `V ≤ 2^8`. (`rounds ≤ 2^14` and `vkLen ≤ 2^10` are structure fields.) -/
def ShapeAdmissible (sz : Sizes) (inWords : ℕ) : Prop :=
  (∀ j, (sh j).Admissible sz) ∧ 8 * (inWords + 3 * (summarize sh inWords).nE + (summarize sh inWords).nIdx) ≤
      TAPE_CAP ∧
    (summarize sh inWords).blowup sz ≤ 2 ^ 8

/-- Tape input words of a protocol: header, `VK`, `x`. -/
def Protocol.inWords {sz : Sizes} (P : Protocol sz) : ℕ := 4 + P.vkLen + sz.nx

/-- The protocol's score inputs. -/
def Protocol.summary {sz : Sizes} (P : Protocol sz) : Summary := summarize P.shape P.inWords

/-- Admission of a protocol at its size (shape caps, `idxLen ≤ 2^(k+8)`), and of a claimed cost
vector (`Q ≤ 2^16`). -/
def Protocol.Admissible {sz : Sizes} (P : Protocol sz) (c : CostVec) : Prop :=
  ShapeAdmissible P.shape sz P.inWords ∧ P.idxShape.IdxAdmissible sz ∧ c.qCount ≤ 2 ^ 16
\end{leancode}

\subsection{\texttt{QProg.lean}: T0 query programs}

A free monad over \code{QUERY}, with free local computation, for the Lean-only track.
\emph{Check}: \lean{eval} and \lean{cost} fail exactly where the machine's \code{QUERY} traps
(oracle or leaf out of range), and \lean{cost} adds exactly \lean{CostVec.addQuery}, the machine's
charge.

\begin{leancode}[fontsize=\scriptsize,label={\footnotesize Contract/QProg.lean}]
import RecursionFast.Contract.Machine

/-- A query program with result type `α`. -/
inductive QProg (α : Type) where
  /-- Stop with result `a`. -/
  | done (a : α)
  /-- QUERY leaf `l` of oracle `j`, continue with the leaf (its `2^logLeaf` elements). -/
  | query (j l : ℕ) (k : List K → QProg α)

namespace QProg
variable {α : Type}

/-- Run against oracles `os` (QUERY numbering as on the machine); an out-of-range oracle or leaf
is `none` (the machine traps there). -/
def eval : QProg α → List Oracle → Option α
  | .done a, _ => some a
  | .query j l k, os =>
    match os[j]? with
    | some o => if l < o.shape.leaves then (k (o.leaf l)).eval os else none
    | none => none

/-- The query counters along the executed path, added to `c` (`none` as for `eval`). -/
def cost : QProg α → List Oracle → CostVec → Option CostVec
  | .done _, _, c => some c
  | .query j l k, os, c =>
    match os[j]? with
    | some o => if l < o.shape.leaves then (k (o.leaf l)).cost os (c.addQuery o.shape) else none
    | none => none

end QProg
\end{leancode}

\subsection{\texttt{Accumulation.lean}: the accumulation track (T2)}

Accumulator shapes (plain part on the next step's tape, at most $2^{12}$ elements; at most 16
oracles of declared shapes), typed accumulator statements (A4-F1), the three reductions
\lean{embed}, \lean{fold}, \lean{decide} with their views, output agreement and cost clauses, and
\lean{Accumulation.Secure} and \lean{Admissible}. \emph{Check}: each view is what the machine sees
(\lean{embed}: \lean{sizes} $\|$ \lean{VK i} $\|$ $x$ with oracle 0 \lean{IDX i}; \lean{fold}: both
plain parts and both accumulators' oracles; \lean{decide}: one); the output oracles are a declared
\lean{select}ion among the step's QUERY oracles, so an output never carries anything the step did
not receive or commit to; \lean{SelectOk} makes the selection well-shaped; \lean{Agrees} includes
equality of the plain output with the words at \lean{OUT\_BASE}; \lean{fold} is secure for two
independent instances (\lean{pairRel}); \lean{decide} outputs the empty accumulator.

\begin{leancode}[fontsize=\scriptsize,label={\footnotesize Contract/Accumulation.lean}]
import RecursionFast.Contract.Cost

open scoped NNReal

/-- Cap on the plain part of an accumulator, in K-elements. -/
def accCap : ℕ := 2 ^ 12
/-- Cap on the number of accumulator oracles. -/
def accOracleCap : ℕ := 16

/-- Declared accumulator shape. -/
structure AccShape where
  plain : ℕ
  oracles : List OShape

/-- An accumulator statement: plain part and one oracle of each declared shape. -/
abbrev AccStmt (A : AccShape) := (Fin A.plain → K) × ((j : Fin A.oracles.length) → OVal (A.oracles.get j))

/-- The accumulator's oracles as the machine sees them (QUERY numbering in declared order). -/
def AccStmt.oracles {A : AccShape} (s : AccStmt A) : List Oracle := List.ofFn fun j => ⟨A.oracles.get j, s.2 j⟩

/-- What a program sees of its input: tape words and QUERY oracles. -/
abbrev View := List Word × List Oracle

/-- Shapes of the oracles a transcript's prover messages produce, in round order. -/
def msgShapes {n : ℕ} (sh : Fin n → Round) : List OShape :=
  (List.ofFn fun i => match sh i with | .msg o => [o] | .chal _ _ => []).flatten

/-- One step: from inputs `Inp` to an accumulator of shape `B`. -/
structure Reduction (Inp WitIn WitOut : Type) (B : AccShape) where
  n : ℕ
  n_le : n ≤ roundCap
  shape : Fin n → Round
  prover : Strategy shape Inp WitIn (AccStmt B) WitOut
  /-- The output plain part, or `none` = reject. -/
  verify : Inp → FullTranscript shape → Option (Fin B.plain → K)
  /-- Output oracle `j` is QUERY oracle `select[j]` of this step (inputs first, then messages). -/
  select : List ℕ
  WitMid : Fin (n + 1) → Type
  extractor : Extractor shape Inp WitIn WitOut WitMid
  ε : ChallengeIdx shape → ℝ≥0

namespace Reduction
variable {Inp WitIn WitOut : Type} {B : AccShape} (R : Reduction Inp WitIn WitOut B) (view : Inp → View)

/-- The step's QUERY oracles: the input's, then the prover messages. -/
def oraclesOf (s : Inp) (tr : FullTranscript R.shape) : List Oracle := (view s).2 ++ tr.oracles

/-- The spec verifier with its full output statement: plain part, and as output oracle `j`
the contents of QUERY oracle `select[j]` (read at the declared length). -/
def verifier : PureVerifier R.shape Inp (AccStmt B) :=
  fun s tr => (R.verify s tr).map fun p =>
    (p, fun j q => ((R.oraclesOf view s tr).getD (R.select.getD j 0) default).get q)

/-- Security of one step for input/output relations. -/
def Secure (relIn : Set (Inp × WitIn)) (relOut : Set (AccStmt B × WitOut)) : Prop :=
  SecureWith relIn relOut R.prover (R.verifier view) R.extractor R.ε

/-- Static check: given input oracle shapes `ins`, the selection yields the declared shapes. -/
def SelectOk (ins : List OShape) : Prop := R.select.map (fun q => (ins ++ msgShapes R.shape)[q]?) =
    B.oracles.map some

/-- The machine run of one step on input `s` and transcript `tr`. -/
def runOn (img : Image) (s : Inp) (tr : FullTranscript R.shape) : Option (CostVec × RiscvZkvm.Rv64.MachineState)
    :=
  run img ((view s).1 ++ tr.tapeWords) (R.oraclesOf view s tr)

/-- **Agreement with output**: on every input and transcript the image accepts iff the spec
does, and then leaves the spec's plain output at `OUT_BASE`. -/
def Agrees (img : Image) : Prop :=
  ∀ s tr, (R.runOn view img s tr).map (fun r => readOut r.2 B.plain) = (R.verify s tr).map wordsOf

/-- Accepting-run cost bound for one step. -/
def AcceptingCostAtMost (img : Image) (c : CostVec) : Prop :=
  ∀ s tr v st, R.runOn view img s tr = some (v, st) → v ≤ c

end Reduction

/-- Two independent instances of a relation. -/
def pairRel {S W : Type} (R : Set (S × W)) : Set ((S × S) × (W × W)) :=
  {p | (p.1.1, p.2.1) ∈ R ∧ (p.1.2, p.2.2) ∈ R}

/-- An accumulation scheme for size `sz`. -/
structure Accumulation (sz : Sizes) where
  acc : AccShape
  acc_le : acc.plain ≤ accCap ∧ acc.oracles.length ≤ accOracleCap
  /-- Holographic preprocessing for `embed`, as in `Protocol`. -/
  vkLen : ℕ
  vkLen_le : vkLen ≤ vkCap
  VK : Index sz → (Fin vkLen → K)
  idxShape : OShape
  IDX : Index sz → OVal idxShape
  /-- The accumulator relation (submitter-chosen). -/
  AccWit : Type
  Racc : Set (AccStmt acc × AccWit)
  embed : Reduction (Index sz × PubInput sz) (Wit sz) AccWit acc
  fold : Reduction (AccStmt acc × AccStmt acc) (AccWit × AccWit) AccWit acc
  decide : Reduction (AccStmt acc) AccWit Unit ⟨0, []⟩

namespace Accumulation
variable {sz : Sizes} (A : Accumulation sz)

/-- `embed` sees `sizes ‖ VK i ‖ x` and oracle 0 = `IDX i`. -/
def embedView : Index sz × PubInput sz → View :=
  fun s => (sz.header ++ wordsOf (A.VK s.1) ++ wordsOf s.2, [⟨A.idxShape, A.IDX s.1⟩])
/-- `fold` sees `sizes ‖ plain₁ ‖ plain₂` and both accumulators' oracles. -/
def foldView : AccStmt A.acc × AccStmt A.acc → View :=
  fun s => (sz.header ++ wordsOf s.1.1 ++ wordsOf s.2.1, AccStmt.oracles s.1 ++ AccStmt.oracles s.2)
/-- `decide` sees `sizes ‖ plain` and the accumulator's oracles. -/
def decideView : AccStmt A.acc → View := fun s => (sz.header ++ wordsOf s.1, AccStmt.oracles s)

/-- **Security of the scheme**: each step secure for its relations; selections well-shaped. -/
structure Secure : Prop where
  embed : A.embed.Secure A.embedView (rel sz) A.Racc
  fold : A.fold.Secure A.foldView (pairRel A.Racc) A.Racc
  decide : A.decide.Secure A.decideView A.Racc Set.univ
  select : A.embed.SelectOk [A.idxShape] ∧ A.fold.SelectOk (A.acc.oracles ++ A.acc.oracles) ∧
    A.decide.SelectOk A.acc.oracles

/-- Shape caps of the three steps (input tape words: header + VK + x, 2·plain, plain). -/
def Admissible : Prop :=
  ShapeAdmissible A.embed.shape sz (4 + A.vkLen + sz.nx) ∧
    ShapeAdmissible A.fold.shape sz (4 + 2 * A.acc.plain) ∧
    ShapeAdmissible A.decide.shape sz (4 + A.acc.plain) ∧ A.idxShape.IdxAdmissible sz ∧
    ∀ o ∈ A.acc.oracles, o.Admissible sz

/-- Score inputs of the three steps (embed, fold, decide). -/
def summary : Summary × Summary × Summary :=
  (summarize A.embed.shape (4 + A.vkLen + sz.nx), summarize A.fold.shape (4 + 2 * A.acc.plain),
    summarize A.decide.shape (4 + A.acc.plain))

end Accumulation
\end{leancode}

\subsection{\texttt{Certificate.lean}: what a record certifies}

One submission structure, one rendered-data structure and one certificate per track, plus the
latency record. \emph{Check}: every field of \lean{Declared0/1/2} is organiser data rendered into
the challenge, and the stub leaves only \lean{submission} open; every clause holds for each of the
three classes; the \lean{summary} and \lean{score} clauses bind the scorer's inputs to the proved
protocols; the T0 cost clause bounds every accepting path of the query program; the T2 admission
caps the query count of each step; there is no decider cap (A5-M8).

\begin{leancode}[fontsize=\scriptsize,label={\footnotesize Contract/Certificate.lean}]
import RecursionFast.Contract.QProg
import RecursionFast.Contract.Accumulation

/-- Image admission: code + data `≤ 2^20` bytes. -/
def imageCap : ℕ := 2 ^ 20

/-! ### T0: Lean-only, query programs -/

/-- A T0 submission: a protocol per class and a query program reading `VK i`, `x` and the
challenges as values and the oracles only through QUERY. -/
structure Submission0 where
  protocol : (c : Fin 3) → Protocol (classSz c)
  qprog : (c : Fin 3) → (Fin (protocol c).vkLen → K) → PubInput (classSz c) →
    Challenges (protocol c).shape → QProg Bool

/-- Rendered T0 data: shape summaries, claimed query counters (all other counters `0`: local
computation is free in T0), claimed `scoreZk4` (the T0 score is the recursion-board formula on
these counters: virtual query cost, FS surcharge and bytes). -/
structure Declared0 where
  summary : Fin 3 → Summary
  claim : Fin 3 → CostVec
  zk4 : ℚ

/-- T0 certificate. -/
structure Certificate0 (S : Submission0) (D : Declared0) : Prop where
  summary : ∀ c, (S.protocol c).summary = D.summary c
  score : scoreZk4 D.summary D.claim = D.zk4
  secure : ∀ c, (S.protocol c).Secure
  admissible : ∀ c, (S.protocol c).Admissible (D.claim c)
  /-- The query program accepts exactly when the spec verifier does (all transcripts). -/
  agrees : ∀ c i x tr, ((S.qprog c ((S.protocol c).VK i) x (fun r => tr r.1)).eval
      ((S.protocol c).oracles i tr) = some true) = ((S.protocol c).verify ((S.protocol c).VK i) x ((S.protocol
          c).IDX i) tr = true)
  /-- Query counters of every accepting path are `≤ claim c`. -/
  cost : ∀ c i x tr v, (S.qprog c ((S.protocol c).VK i) x (fun r => tr r.1)).eval ((S.protocol c).oracles i tr)
      = some true →
      (S.qprog c ((S.protocol c).VK i) x (fun r => tr r.1)).cost ((S.protocol c).oracles i tr) {} = some v → v ≤
          D.claim c

/-! ### T1: RISC-V verifier -/

/-- A T1 submission: a protocol per class (the image is rendered challenge data). -/
structure Submission1 where
  protocol : (c : Fin 3) → Protocol (classSz c)

/-- Rendered T1 data: the image, shape summaries, claimed counters, claimed `scoreZk4`. -/
structure Declared1 where
  image : Image
  summary : Fin 3 → Summary
  claim : Fin 3 → CostVec
  zk4 : ℚ

/-- T1 certificate (DESIGN §4, §5, §15). -/
structure Certificate1 (S : Submission1) (D : Declared1) : Prop where
  summary : ∀ c, (S.protocol c).summary = D.summary c
  score : scoreZk4 D.summary D.claim = D.zk4
  secure : ∀ c, (S.protocol c).Secure
  admissible : ∀ c, (S.protocol c).Admissible (D.claim c)
  imageSize : D.image.byteSize ≤ imageCap
  agrees : ∀ c, (S.protocol c).Agrees D.image
  cost : ∀ c, (S.protocol c).AcceptingCostAtMost D.image (D.claim c)

/-- Latency-board record (on top of `Certificate1`): claimed `(W, S)` per class, `scoreLat3`. -/
structure Latency1 (S : Submission1) (D : Declared1) (lat : Fin 3 → ℕ × ℕ) (lat3 : ℚ) : Prop where
  bound : ∀ c, (S.protocol c).AcceptingLatencyAtMost D.image (lat c).1 (lat c).2
  score : scoreLat3 lat = lat3

/-! ### T2: accumulation -/

/-- A T2 submission: an accumulation scheme per class (the three images are challenge data). -/
structure Submission2 where
  scheme : (c : Fin 3) → Accumulation (classSz c)

/-- Claimed cost vectors of the three steps at one class. -/
structure Claim2 where
  embed : CostVec
  fold : CostVec
  decide : CostVec

/-- `S⁴` of (embed, fold, decide) per profile and class. The T2 score is
`Score_T2 = (Π_P Π_k (a^(1/4) + b^(1/4) + c^(1/4)/64))^(1/12)` with `(a, b, c)` these values
(DESIGN §6: `S_leaf + S_fold + S_dec/64`, geomean over profiles and classes). -/
def scoreT2Fourths (d : Fin 3 → Summary × Summary × Summary) (c : Fin 3 → Claim2) (P : Fin 4) (k : Fin 3) :
    ℚ × ℚ × ℚ :=
  ((d k).1.scoreFourth (profiles P) (classSz k) (c k).embed,
    (d k).2.1.scoreFourth (profiles P) (classSz k) (c k).fold,
    (d k).2.2.scoreFourth (profiles P) (classSz k) (c k).decide)

/-- Rendered T2 data: the three images, summaries, claims, claimed `S⁴` table. -/
structure Declared2 where
  embedImage : Image
  foldImage : Image
  decideImage : Image
  summary : Fin 3 → Summary × Summary × Summary
  claim : Fin 3 → Claim2
  fourths : Fin 4 → Fin 3 → ℚ × ℚ × ℚ

/-- T2 certificate. There is no decider cap (DESIGN §15.3, A5-M8, superseding §7's
`C_dec ≤ 2^32`): the decider is bounded only by the fuel `2^36` and enters the score at `/64`, so
linear (Arc/WARP-style) deciders are admissible. -/
structure Certificate2 (S : Submission2) (D : Declared2) : Prop where
  summary : ∀ c, (S.scheme c).summary = D.summary c
  score : ∀ P k, scoreT2Fourths D.summary D.claim P k = D.fourths P k
  secure : ∀ c, (S.scheme c).Secure
  admissible : ∀ c, (S.scheme c).Admissible ∧ (D.claim c).embed.qCount ≤ 2 ^ 16 ∧
    (D.claim c).fold.qCount ≤ 2 ^ 16 ∧ (D.claim c).decide.qCount ≤ 2 ^ 16
  imageSize : D.embedImage.byteSize ≤ imageCap ∧ D.foldImage.byteSize ≤ imageCap ∧
    D.decideImage.byteSize ≤ imageCap
  agrees : ∀ c, (S.scheme c).embed.Agrees (S.scheme c).embedView D.embedImage ∧
    (S.scheme c).fold.Agrees (S.scheme c).foldView D.foldImage ∧
    (S.scheme c).decide.Agrees (S.scheme c).decideView D.decideImage
  cost : ∀ c, (S.scheme c).embed.AcceptingCostAtMost (S.scheme c).embedView D.embedImage (D.claim c).embed ∧
    (S.scheme c).fold.AcceptingCostAtMost (S.scheme c).foldView D.foldImage (D.claim c).fold ∧
    (S.scheme c).decide.AcceptingCostAtMost (S.scheme c).decideView D.decideImage (D.claim c).decide
\end{leancode}

\subsection{Auditor's checklist}\label{subsec:app-checklist}

The checklist merges the definitions review (A2, \enquote{what a human auditor must check}) and the
contract review's auditability assessment (A4), updated for v0.3. Statement-level holes are the one
class of error no checker detects (\cref{subsec:bg-lean}), so this reading is the last line of
defense.

\paragraph{Lean text, in reading order.}
\begin{enumerate}[nosep,leftmargin=1.6em]
\item \lean{K}, \lean{E} and their encodings (\code{Field.lean}); the library controls \lean{K\_spec},
  \lean{card\_K}, \lean{card\_E}.
\item \lean{R1CS}, $z=(1,x,w)$, \lean{rel}, the sparse \lean{Index} and its \lean{nnz\_le}, the class
  shapes; padding and \lean{pad\_sat\_iff}.
\item \lean{Round}, \lean{Round.Ty}: message types are raw \lean{OVal o}, challenge types the
  contract's $E^{n_E}\times\prod_b\{0,1\}^b$; \lean{Round.tapeWords}.
\item \lean{Strategy.transcript}, \lean{PerfectComplete}; \lean{KSF} (three laws), \lean{Bad},
  \lean{RbrKS} (counting form), read as backward extraction, worst case over prefixes;
  \lean{SecureWith} with \lean{cap}.
\item \lean{Protocol} caps (\lean{n ≤ roundCap}, \lean{vkLen ≤ vkCap}) and \lean{Protocol.Secure}.
\item The machine: decoder arms, memory checks, system calls, \code{QUERY} closed form, fuel,
  \lean{HALT}; the hart discipline; that the accounting and the timer never drive the decision; the
  cost classes (\lean{kind}) and first-touch keys.
\item \lean{Agrees}: quantified over all indices, inputs and raw transcripts; tape layout and oracle
  numbering.
\item Cost clauses and score functions: componentwise order, accepting-run quantifier, profile
  weights against the published table, \lean{Summary}, \lean{fs}, \lean{bytes}, \lean{blowup},
  \lean{scoreFourth}, \lean{scoreZk4}, \lean{scoreLat3}; admission predicates.
\item \lean{QProg.eval} and \lean{QProg.cost} against \code{QUERY}.
\item The accumulation shape, views, \lean{select}/\lean{SelectOk}, the three \lean{Reduction.Secure}
  and output agreement.
\item Each \lean{Certificate*} clause and the challenge stubs: only \lean{submission} is open.
\end{enumerate}

\paragraph{Outside the Lean text.}
\begin{enumerate}[nosep,leftmargin=1.6em,resume]
\item Harness prose and scripts (\code{harness.sh}, \code{Gate.lean}, \code{olean-manifest.py},
  \code{render-challenge.py}): the renderer and claim schema with path confinement; the stage order of
  Decision~A-3 (comparator \emph{first}, in a fresh tree in which no solution code has run, with the
  organiser package read-only; then the gate in a second fresh tree, whose build must be
  byte-identical to the comparator-verified one (round~4, A7-M3); the badge as its own process);
  the \code{.olean} manifest checked before and after each stage; the environment gate's trust rule (a
  module is trusted iff every \code{.olean} part matches the organiser manifest by content hash, never
  by name, location or import graph: \lean{SpoofCert}, A6-G2) and its rejection list (non-standard
  axioms, \lean{sorryAx}, \lean{Lean.ofReduceBool}; \lean{implemented\_by}, \lean{extern},
  \lean{csimp}, \lean{init}, \lean{builtin\_init} and \lean{export} entries in any untrusted module,
  the organiser library included; rule~4 on the badge-path fields only); the axiom audit; two kernel
  replays (Lean kernel and nanoda);
  the production sandbox.
\item Fiat--Shamir prefix, sampler and canonical encodings (\cref{subsec:sec-fs}); the canonical
  form of benchmark indices. None of these is implemented yet (H-3).
\item The emulator against \lean{run}. No organiser emulator exists yet (H-4); the baseline's
  unverified \code{rfemu} was differential-tested against \lean{runFull} on two tiny instances
  (counters, first touch and timer; harts not covered), and the QEMU tests of A4 cover register, M,
  Zbc, W-op and load/store semantics but not control flow or fetch.
\item The cited meta-theorems with exact numbers: Chiesa--Yogev Thms.~31.2.1 and 32.8.9 (soundness),
  WARP Thm.~B.4 with Chiesa--Yogev Thm.~32.8.10 (knowledge, extraction time unbounded at L1), and the
  side conditions $t\ge t^*$ of \cref{subsec:sec-cap}.
\end{enumerate}

\paragraph{Auditability.} \emph{Size}: 796 contract code lines plus the vendored code they reach
(474 declarations, 930 source lines: \lean{execInstrBr}, \lean{decode}, \lean{MachineState}, and
ZisK accelerator code that is dead at run time), and the meaning of a few core and Mathlib notions:
\lean{Nat.card}, \lean{Finset.card} with \lean{Ico} and \lean{image}, \lean{Set}, \lean{ℝ≥0},
\lean{ℚ} with ceiling, \lean{BitVec}, \lean{List.ofFn}, \lean{![…]} (\lean{Matrix.vecCons}, used by
\lean{profiles}), \lean{Fin.foldl}, \lean{Fin.cons}, \lean{Fin.append}, \lean{Fin.snoc},
\lean{Fin.induction}, \lean{List.Forall₂}. Beyond the text, the trusted surface includes the Lean
elaborator (for the challenge), Lake, the Lean compiler (gate, axiom audit, badge path), the
comparator, \code{lean4export} and nanoda binaries (not covered by the \code{.olean} manifest), and
the host tools of the harness (Python~$\ge$3.11, bubblewrap, \code{sha256sum}); see G-7.
\emph{Time}: A4 estimated, for v0.2, 1--1.5 days for the contract and one more day for the machine and
vendored semantics for a cryptographer who reads Lean, and 3--5 days without Lean experience; v0.3
is 37\% larger, almost all of it in \code{Machine.lean} and \code{Cost.lean}, so these estimates
grow by roughly a day (our extrapolation, not a measurement). \emph{Hardest definitions}, in order:
(1)~\code{Security.lean} (prefix indexing by \lean{Fin (n+1)}, \lean{castSucc}/\lean{succ},
\lean{cast E.eqIn}, the $\exists w$ in \lean{Bad}); (2)~\code{Machine.lean}, now with harts: the
nested \lean{execute} at \code{SPAWN}, the \lean{allows} discipline, the threaded accounting, and
memory as \lean{Word → Word} keyed by aligned address (defined on every key, read only at aligned
ones); (3)~\code{Accumulation.lean} (views, \lean{select} with \lean{getD default}; the untyped
oracle lists where A4-F1 hid are now typed by \lean{AccStmt}); (4)~\code{Cost.lean} (score
functions, now wired into the certificates by the \lean{summary} and \lean{score} clauses, which
closes A4's F2 and F4 inside the statement). \emph{What remains outside Lean}: the emulator, the
Fiat--Shamir derivation and sampler, the canonical index form, the badge checks, the submission
preflight and source policy, and the production sandbox, most of which are specified but not yet
built (H-9); an auditor cannot sign off on deployment meaning until these are reviewed with
their test vectors.
\endgroup
